%% Question bank chapter file - generated by bank_tools.py build (do not edit while the bank is still being built:
%% it is regenerated from the per-test data in Question Bank\data). Every question is one bankq block:
%%   \begin{bankq}[id=..., source={...}, paper=.., qno=.., type=..., subtopic={...}, answer={...}, answer_src=key|solved,
%%                 working={...}, difficulty=1-10, diff_parts={K S I C F}, diff_note={...},
%%                 group=<passage id or empty>, also_in={duplicate copies in other tests}]
%% difficulty = min(10, K+S+I+C+F-1): K concepts 1-3, S steps 1-4, I insight 0-2, C computation 0-2, F format 0-1
%%   ...question LaTeX, options, figures (\bfig / \bfigi -> img/)...
%%   \end{bankq}
%% Passages shared by several questions are bankpassage blocks; their questions carry group=<passage id>.
\documentclass[10pt]{article}
\usepackage[a4paper,margin=1.8cm]{geometry}
\usepackage[T1]{fontenc}
\usepackage{amsmath,amssymb,mathtools}
\usepackage{graphicx}
\usepackage[export]{adjustbox}
\usepackage{xcolor,needspace,keyval}
\usepackage{enumitem}
\usepackage{booktabs,array,tabularx}
\usepackage[hidelinks]{hyperref}
\setlength{\parindent}{0pt}
\setlength{\parskip}{3pt}
\makeatletter
\define@key{bq}{id}{\def\bq@id{#1}}\define@key{bq}{source}{\def\bq@src{#1}}\define@key{bq}{paper}{\def\bq@paper{#1}}
\define@key{bq}{qno}{\def\bq@qno{#1}}\define@key{bq}{type}{\def\bq@type{#1}}\define@key{bq}{subtopic}{\def\bq@sub{#1}}
\define@key{bq}{answer}{\def\bq@ans{#1}}\define@key{bq}{answer_src}{\def\bq@asrc{#1}}\define@key{bq}{working}{\def\bq@work{#1}}
\define@key{bq}{group}{\def\bq@group{#1}}\define@key{bq}{also_in}{\def\bq@also{#1}}
\define@key{bq}{difficulty}{\def\bq@diff{#1}}\define@key{bq}{diff_parts}{\def\bq@dparts{#1}}\define@key{bq}{diff_note}{\def\bq@dnote{#1}}
\newenvironment{bankq}[1][]{\def\bq@work{}\def\bq@also{}\def\bq@group{}\def\bq@diff{?}\def\bq@dparts{}\def\bq@dnote{}\setkeys{bq}{#1}%
  \par\needspace{4\baselineskip}\medskip\noindent{\small\textbf{\bq@id}\hfill\color{gray}[\bq@src, Paper \bq@paper, Q\bq@qno]
  \bq@type\ $\cdot$ \bq@sub\ $\cdot$ \textbf{D\bq@diff}}\par\nopagebreak\relax}
  {\par{\small\color{teal}Answer: \textbf{\bq@ans}\ifx\bq@work\@empty\else\ (solved: \bq@work)\fi
  \ifx\bq@also\@empty\else\par\color{gray}Also in: \bq@also\fi}\par}
\define@key{bp}{id}{\def\bp@id{#1}}\define@key{bp}{source}{\def\bp@src{#1}}\define@key{bp}{questions}{\def\bp@qs{#1}}
\newenvironment{bankpassage}[1][]{\setkeys{bp}{#1}\par\needspace{4\baselineskip}\bigskip\noindent
  {\small\textbf{Passage \bp@id}\hfill\color{gray}[\bp@src; questions \bp@qs]}\par\nopagebreak\relax\itshape}{\par}
\makeatother
\newcommand{\bfig}[2]{\par\smallskip{\centering\includegraphics[width=#1pt,max width=0.62\linewidth,max height=0.32\textheight]{img/#2}\par}\smallskip}
\newcommand{\bfigi}[2]{\includegraphics[width=#1pt,max width=\linewidth,max height=0.32\textheight,valign=c]{img/#2}}
\newcommand{\bmissing}[1]{\textit{[figure missing in the source conversion: \detokenize{#1}]}}
\newenvironment{qoptions}{\begin{itemize}[leftmargin=1.8em,itemsep=1pt,topsep=2pt,label={}]}{\end{itemize}}
\newcommand{\opt}[2]{\item[(#1)] #2}
\providecommand{\fwstart}{}\providecommand{\fwend}{}\providecommand{\qfullsplit}{}
\newcolumntype{L}[1]{>{\raggedright\arraybackslash}p{#1}}
\newcolumntype{C}[1]{>{\centering\arraybackslash}p{#1}}
\newcolumntype{R}[1]{>{\raggedleft\arraybackslash}p{#1}}
\providecommand{\sectionhdr}[1]{}\providecommand{\subsectionhdr}[1]{}


\begin{document}
{\Large\bfseries Electrostatics --- Question Bank}\par\smallskip
{\small Physics chapter P08. Questions from all converted test papers, grouped by source series. Answers: printed key where it matches the paper, otherwise solved (answer\_src=solved, with working). Exact repeats across papers appear once, with also\_in listing the other copies.}\par\medskip\hrule

\section*{TOAS papers}
\begin{bankq}[id=TOAS04-P1-PHY-S4-Q08, source={TOAS - Test 04 (15 Mar 2022)}, paper=1, qno=8, type=numerical, subtopic={Coulomb's law}, answer={225}, answer_src=key, difficulty=4, diff_parts={K2 S1 I1 C0 F1}, diff_note={concepts: Coulomb force, motion in a uniform field | steps: equal accelerations condition; Coulomb force between the charges and distance | I: constant separation means equal accelerations}]
Two charged particles $(m,q)$ and $(2m,-2q)$ are placed in a gravity free space where a uniform electric field $E$ exists as shown in figure. After the particles are released, they stay at a constant distance from each other. What is the distance between them in centimeter? (Neglected the gravitational interaction). ($K$ is electrostatic constant and take $\dfrac{Kq}{E}=\dfrac{27}{8}$)
\bfig{101.8}{TOAS04_p1_phy_s4_q08_1.png}
\end{bankq}
\begin{bankq}[id=TOAS04-P2-PHY-S1-Q04, source={TOAS - Test 04 (15 Mar 2022)}, paper=2, qno=4, type=multi, subtopic={Electric field}, answer={A,C,D}, answer_src=key, difficulty=4, diff_parts={K2 S2 I0 C0 F1}, diff_note={concepts: field of a sheet, a point charge | steps: field of the sheet is sigma/2 epsilon0; zero-field point from the graph gives Q; superposition at a/2 and 5a/2}]
The X-Y-Z space has an infinite charged sheet having surface charge density $\sigma$ on the Y-Z plane. A negative charge particle $-Q$ is placed at $(a,0,0)$. The variation of net electric field on $x$-axis is shown in the figure. Then choose the correct option(s) :
\bfig{120.1}{TOAS04_p2_phy_s1_q04_1.png}
\begin{qoptions}
\opt{A}{$Q=\sigma\,2\pi a^2$}
\opt{B}{$Q=\sigma\,4\pi a^2$}
\opt{C}{The net electric field at the point $\left(\dfrac a2,0,0\right)$ is $\dfrac{5\sigma}{2\varepsilon_0}$}
\opt{D}{At point $\left(\dfrac{5a}2,0,0\right)$ the net electric field is $\dfrac{5\sigma}{18\varepsilon_0}$.}
\end{qoptions}
\end{bankq}
\begin{bankq}[id=TOAS04-P2-PHY-S1-Q05, source={TOAS - Test 04 (15 Mar 2022)}, paper=2, qno=5, type=multi, subtopic={Gauss's law}, answer={A,C}, answer_src=key, difficulty=6, diff_parts={K2 S2 I1 C1 F1}, diff_note={concepts: Gauss's law, charge by integration | steps: total charge by integrating rho; enclosed charge within radius r; field from Gauss's law; potential | I: enclosed charge by integration}]
A point charge $q$ is kept at origin. The space around origin can be considered as non conducting sphere, tends to infinity and having continuous charge distribution, whose volume charge density is
\[\rho=\frac{-q}{4\pi a}\cdot\frac{e^{-r/a}}{r^2}\]
Where `$a$' is a positive constant and `$r$' is distance of a point from origin. The potential at infinity can be considered to be zero. The electric constant of space is 1. (Neglect the effect of polarization)
\begin{qoptions}
\opt{A}{The total charge of system is zero.}
\opt{B}{The total charge of system is $\dfrac q2$.}
\opt{C}{The magnitude of electrostatic field at a distance of $r$ from origin is $\dfrac{qe^{-r/a}}{4\pi\varepsilon_0r^2}$}
\opt{D}{The magnitude of electrostatic field at a distance of $r$ from origin is $\dfrac{qe^{-r/a}}{8\pi\varepsilon_0r^2}$}
\end{qoptions}
\end{bankq}
\begin{bankq}[id=TOAS05-P1-PHY-S1-Q05, source={TOAS - Test 05 (03 Apr 2022)}, paper=1, qno=5, type=multi, subtopic={Electric field}, answer={B,D}, answer_src=key, difficulty=2, diff_parts={K1 S1 I0 C0 F1}, diff_note={concepts: electric field of a nonuniform ring | steps: field of an element; integrate x and y components for lambda0 cos theta}]
A thin non conducting ring of radius $R$ has linear charge density $\lambda$ varying as $\lambda=\lambda_0\cos\theta$, where $\lambda_0$ is a positive constant and $\theta$ is angle measured from positive $x$-direction as shown. Consider centre of ring is at origin.
\bfig{130.1}{TOAS05_p1_phy_s1_q05_1.png}
\begin{qoptions}
\opt{A}{Resultant electric field at origin is towards positive $y$-axis.}
\opt{B}{Resultant electric field at origin is towards negative $x$-axis.}
\opt{C}{Magnitude of resultant electric field at origin is $\dfrac{\lambda_0}{8\varepsilon_0R}$.}
\opt{D}{Magnitude of resultant electric field at origin is $\dfrac{\lambda_0}{4\varepsilon_0R}$.}
\end{qoptions}
\end{bankq}
\begin{bankq}[id=TOAS05-P1-PHY-S1-Q06, source={TOAS - Test 05 (03 Apr 2022)}, paper=1, qno=6, type=multi, subtopic={Electric field}, answer={A,C}, answer_src=key, difficulty=2, diff_parts={K1 S1 I0 C0 F1}, diff_note={concepts: block on an incline in a horizontal field | steps: normal force from the components; zero acceleration condition}]
A fixed smooth inclined plane making an angle $\theta$ with the horizontal is placed in a uniform horizontal electric field $E_0$ as shown in figure. A small block of mass $m$ having charge $+q$ is released on inclined plane. (Acceleration due to gravity is `$g$')
\bfig{95.0}{TOAS05_p1_phy_s1_q06_1.png}
\begin{qoptions}
\opt{A}{Normal contact force applied by inclined plane on block is $(mg\cos\theta+qE_0\sin\theta)$}
\opt{B}{Normal contact force force applied by inclined plane on block is $(mg\cos\theta-qE_0\sin\theta)$}
\opt{C}{If acceleration of block is zero then $\theta=\tan^{-1}\left(\dfrac{qE_0}{mg}\right)$.}
\opt{D}{If acceleration of block is zero then $\theta=\tan^{-1}\left(\dfrac{mg}{qE_0}\right)$.}
\end{qoptions}
\end{bankq}
\begin{bankq}[id=TOAS05-P1-PHY-S3-Q05, source={TOAS - Test 05 (03 Apr 2022)}, paper=1, qno=5, type=numerical, subtopic={Electric field}, answer={3.50}, answer_src=key, difficulty=3, diff_parts={K1 S2 I0 C0 F1}, diff_note={concepts: equilibrium on the axis of four charges | steps: vertical force from the four charges; mg equals that force; compare with the given form}]
Four small particles charged with equal positive charge $+Q$ each are arranged at the four corners of a horizontal square of side $a$. A particle having charge $+q$ and mass $m$ is placed at a height $h$ above the centre of the square. Condition under which particle at height $h$ is in equilibrium is given by
\[m=\frac{Qqh}{g\pi\varepsilon_0\left(h^2+\frac{a^2}{N_1}\right)^{N_2}}.\]
Fill the value of $(N_1+N_2)$. (Acceleration due to gravity is `$g$')
\bfig{140.6}{TOAS05_p1_phy_s3_q05_1.png}
\end{bankq}
\begin{bankq}[id=TOAS05-P2-PHY-S1-Q01, source={TOAS - Test 05 (03 Apr 2022)}, paper=2, qno=1, type=single, subtopic={Gauss's law}, answer={A}, answer_src=key, difficulty=1, diff_parts={K1 S1 I0 C0 F0}, diff_note={concepts: Gauss's law and nonuniform ball | steps: integrate rho and apply Gauss's law}]
A ball of radius $R$ carries a positive charge whose volume charge density depends only on separation $r$ from the ball's centre as $\rho=\rho_0\left(1-\dfrac rR\right)$ where $\rho_0$ is constant. The electric field at distance $r_1$ from centre (For $r_1<R$) is :-
\begin{qoptions}
\opt{A}{$\dfrac{\rho_0r_1}{3\epsilon_0}\left(1-\dfrac{3r_1}{4R}\right)$}
\opt{B}{$\dfrac{\rho_0r_1}{3\epsilon_0}\left(1-\dfrac{r_1}{R}\right)$}
\opt{C}{$\dfrac{\rho_0r_1}{2\epsilon_0}\left(1-\dfrac{r_1}{2R}\right)$}
\opt{D}{None of these}
\end{qoptions}
\end{bankq}
\begin{bankq}[id=TOAS05-P2-PHY-S1-Q02, source={TOAS - Test 05 (03 Apr 2022)}, paper=2, qno=2, type=single, subtopic={Coulomb's law}, answer={A}, answer_src=key, difficulty=1, diff_parts={K1 S1 I0 C0 F0}, diff_note={concepts: orbital motion under Coulomb force | steps: Coulomb force equals m r omega\^{}2}]
A particle of mass $m$ and carrying charge $-q_1$ is moving around a fixed charge $+q_2$ along a circular path of radius $r$. Find period of revolution of the charge $-q_1$ about $+q_2$.
\begin{qoptions}
\opt{A}{$\sqrt{\dfrac{16\pi^3\epsilon_0mr^3}{q_1q_2}}$}
\opt{B}{$\sqrt{\dfrac{8\pi^3\epsilon_0mr^3}{q_1q_2}}$}
\opt{C}{$\sqrt{\dfrac{q_1q_2}{16\pi^3\epsilon_0mr^3}}$}
\opt{D}{Zero}
\end{qoptions}
\end{bankq}
\begin{bankq}[id=TOAS05-P2-PHY-S1-Q03, source={TOAS - Test 05 (03 Apr 2022)}, paper=2, qno=3, type=single, subtopic={Gauss's law}, answer={A}, answer_src=key, difficulty=1, diff_parts={K1 S1 I0 C0 F0}, diff_note={concepts: Gauss's law for a given field | steps: net flux through the cube equals the charge over epsilon0}]
In a region, electric field depends on $X$-axis as $E=E_0x^2$. There is a cube of edge $a$ as shown. Then find the charge enclosed in that cube.
\bfig{128.0}{TOAS05_p2_phy_s1_q03_1.png}
\begin{qoptions}
\opt{A}{$5\epsilon_0a^4E_0$}
\opt{B}{$3\epsilon_0a^4E_0$}
\opt{C}{$4\epsilon_0a^4E_0$}
\opt{D}{Zero}
\end{qoptions}
\end{bankq}
\begin{bankq}[id=TOAS05-P2-PHY-S1-Q05, source={TOAS - Test 05 (03 Apr 2022)}, paper=2, qno=5, type=single, subtopic={Electric field}, answer={C}, answer_src=key, difficulty=2, diff_parts={K1 S1 I1 C0 F0}, diff_note={concepts: field inside a cavity | steps: uniform field inside the cavity so a straight line motion | I: superposition of a negative sphere}]
A spherical cavity is made in a uniformly charged non-conducting sphere and a positive charge is released from position $P$ as shown in figure. Then path followed by particle is :-
\bfig{112.7}{TOAS05_p2_phy_s1_q05_1.png}
\begin{qoptions}
\opt{A}{A straight line along radially inward of cavity}
\opt{B}{A straight line along radially outward of cavity}
\opt{C}{A straight line along positive $X$-direction}
\opt{D}{A straight line along radially outward of sphere}
\end{qoptions}
\end{bankq}
\begin{bankq}[id=TOAS05-P2-PHY-S1-Q16, source={TOAS - Test 05 (03 Apr 2022)}, paper=2, qno=16, type=single, subtopic={Electric field}, answer={B}, answer_src=key, difficulty=2, diff_parts={K1 S1 I1 C0 F0}, diff_note={concepts: electric field of unequal opposite charges | steps: sign of the field in the three regions; position of the null point | I: null point lies beyond the smaller charge}]
A positive point charge `$Q$' is located at origin and a negative point charge `$q$' is located on $x$-axis at a distance `$r$' from `$Q$'. Magnitude of charge of `$Q$' is greater than that of `$q$'. Then, which of the following may represent the $E$-$x$ graph?
\begin{qoptions}
\opt{A}{\bfigi{104.1}{TOAS05_p2_phy_s1_q16_A.png}}
\opt{B}{\bfigi{106.0}{TOAS05_p2_phy_s1_q16_B.png}}
\opt{C}{\bfigi{102.6}{TOAS05_p2_phy_s1_q16_C.png}}
\opt{D}{\bfigi{106.0}{TOAS05_p2_phy_s1_q16_D.png}}
\end{qoptions}
\end{bankq}
\begin{bankq}[id=TOAS05-P2-PHY-S2-Q01, source={TOAS - Test 05 (03 Apr 2022)}, paper=2, qno=1, type=numerical, subtopic={Gauss's law}, answer={16}, answer_src=key, difficulty=2, diff_parts={K1 S1 I0 C0 F1}, diff_note={concepts: Gauss's law and charge density Kr | steps: enclosed charge integral at R/2 and E}]
A sphere of radius $R$ have volume charge density give as $\rho(r)=Kr$ for $r\le R$ and 0 for $r>R$. If electric field at distance $\dfrac R2$ from centre is $\dfrac{KR^2}{N\epsilon_0}$, then $N$ will be :
\end{bankq}
\begin{bankq}[id=TOAS06-P1-PHY-S1-Q01, source={TOAS - Test 06 (17 Apr 2022)}, paper=1, qno=1, type=single, subtopic={Electric potential}, answer={A}, answer_src=key, difficulty=2, diff_parts={K1 S1 I1 C0 F0}, diff_note={concepts: potential of a charged hemisphere | steps: potential at the centre of the base is half that of the full sphere's centre; work to take the charge to infinity | I: half-sphere symmetry argument}]
A solid non-conducting hemisphere of radius $R$ has a uniformly distributed positive charge of volume density $\rho$. A negatively charged particle having charge $q$ is transferred from centre of its base to infinity. Work performed in the process is :
\begin{qoptions}
\opt{A}{$\dfrac{q\rho R^{2}}{4\in_0}$}
\opt{B}{$\dfrac{q\rho R^{2}}{2\in_0}$}
\opt{C}{$\dfrac{q\rho R^{2}}{8\in_0}$}
\opt{D}{$\dfrac{-q\rho R^{2}}{8\in_0}$}
\end{qoptions}
\end{bankq}
\begin{bankq}[id=TOAS06-P1-PHY-S1-Q02, source={TOAS - Test 06 (17 Apr 2022)}, paper=1, qno=2, type=single, subtopic={Electric potential}, answer={B}, answer_src=key, difficulty=1, diff_parts={K1 S1 I0 C0 F0}, diff_note={concepts: field from a potential | steps: E = -grad phi = (-y, -x); direction on each axis and at (2,2)}]
The potential of electrostatic field is given by $\phi=xy$, where $x$ \& $y$ are coordinate. Which of the following is correct ?
\begin{qoptions}
\opt{A}{Electric field at all points on positive $x$-axis is in positive $y$-direction.}
\opt{B}{Electric field at all points of positive $y$-axis is in negative $x$-direction.}
\opt{C}{Electric field at $(2,2)$ is $\sqrt{2}\,$N/m.}
\opt{D}{Electric field at $x,y$ makes an angle $\theta=\tan^{-1}\!\left(\dfrac{y}{x}\right)$ with $x$-axis.}
\end{qoptions}
\end{bankq}
\begin{bankq}[id=TOAS06-P1-PHY-S1-Q03, source={TOAS - Test 06 (17 Apr 2022)}, paper=1, qno=3, type=single, subtopic={Conductors}, answer={C}, answer_src=key, difficulty=3, diff_parts={K1 S2 I1 C0 F0}, diff_note={concepts: concentric shells and earthed inner shell | steps: potential of the innermost shell is zero; relate charges shell by shell; solve for q1 | I: recursive relation between shell charges}]
Consider a system of $n$ concentric conducting shells of radii $r,2r,4r,8r,\ldots..,2^{n-1}r$ respectively. If inner shell is earthed and $Q$ charge is given on outer most shell, charge appears on inner most shell is :-
\bfig{92.3}{TOAS06_p1_phy_s1_q03_1.png}
\begin{qoptions}
\opt{A}{$-Q$}
\opt{B}{$-\dfrac{Q}{2^{n}}$}
\opt{C}{$-\dfrac{Q}{2^{n-1}}$}
\opt{D}{$+Q$}
\end{qoptions}
\end{bankq}
\begin{bankq}[id=TOAS06-P1-PHY-S2-Q06, source={TOAS - Test 06 (17 Apr 2022)}, paper=1, qno=6, type=multi, subtopic={Conductors}, answer={A,C,D}, answer_src=key, difficulty=3, diff_parts={K1 S2 I0 C0 F1}, diff_note={concepts: charge distribution on parallel plates | steps: charge conservation on each isolated pair; earthing condition when the key closes; final charges}]
Initially two plates (A and B) are given charges as shown. Find the final charges on the plates after the key $k$ is closed. All plates are conducting, parallel and of infinite length and breadth.
\bfig{70.8}{TOAS06_p1_phy_s2_q06_1.png}
\begin{qoptions}
\opt{A}{Final charges on plate A is $Q$}
\opt{B}{Final charge on plate B is $0$}
\opt{C}{Final charge on plate C is $3Q/2$}
\opt{D}{Final charge on plate D is $0$}
\end{qoptions}
\end{bankq}
\begin{bankq}[id=TOAS06-P1-PHY-S2-Q11, source={TOAS - Test 06 (17 Apr 2022)}, paper=1, qno=11, type=multi, subtopic={Dipoles}, answer={A,C,D}, answer_src=key, difficulty=3, diff_parts={K1 S2 I0 C0 F1}, diff_note={concepts: dipole in the field of a point charge | steps: net force from field gradient; torque about the centre; magnitude of force; equilibrium of the free system}]
A dipole of dipole moment $\vec{P}$ is placed in front of a point charge as shown. Choose the correct statement(s) :
\bfig{81.3}{TOAS06_p1_phy_s2_q11_1.png}
\begin{qoptions}
\opt{A}{The dipole will experience a net force irrespective of value of $\theta$.}
\opt{B}{The dipole will experience a net torque about its center irrespective of value of $\theta$.}
\opt{C}{Force experience by dipole is given by $\dfrac{kQP}{r^{3}}\sqrt{1+3\cos^{2}\theta}\,$.}
\opt{D}{If charge $Q$ is not fixed, the whole system can't be in equilibrium for any value of $r$ \& $\theta$.}
\end{qoptions}
\end{bankq}
\begin{bankq}[id=TOAS06-P1-PHY-S2-Q12, source={TOAS - Test 06 (17 Apr 2022)}, paper=1, qno=12, type=multi, subtopic={Electric field}, answer={A,D}, answer_src=key, difficulty=5, diff_parts={K2 S2 I1 C0 F1}, diff_note={concepts: field of a semi-infinite wire, tension in a ring | steps: field components at points of the ring; net force on the ring; tension from half ring balance | I: symmetry of horizontal components}]
Diagram shows a semi-infinite wire having a charge density $\lambda$. A ring carrying a total charge $Q$ is kept as shown. Choose the correct option(s) :
\bfig{104.1}{TOAS06_p1_phy_s2_q12_1.png}
\begin{qoptions}
\opt{A}{The force experienced by ring is in vertical direction \& have a value of $\dfrac{\lambda Q}{4\pi\in_0R}$}
\opt{B}{Net force experienced by ring is in vertical direction and have value of $\dfrac{\lambda Q}{2\pi\in_0R}$}
\opt{C}{Tension in the ring is zero.}
\opt{D}{Tension in the ring is $\dfrac{Q\lambda}{8\pi^{2}\in_0R}$}
\end{qoptions}
\end{bankq}
\begin{bankq}[id=TOAS06-P1-PHY-S3-Q02, source={TOAS - Test 06 (17 Apr 2022)}, paper=1, qno=2, type=numerical, subtopic={Electric field}, answer={9.00}, answer_src=key, difficulty=5, diff_parts={K2 S2 I1 C0 F1}, diff_note={concepts: Coulomb force, volume integration | steps: set up spherical coordinates; z-component of the force; integrate r, theta, phi; substitute numbers | I: spherical coordinates and symmetry}]
Find the coulomb attraction (in N) on a charge of magnitude $5\times10^{-4}$C at the origin due to a charge (of constant density $\rho$) occupying the volume inside the cone $z^2=x^2+y^2$, $0<z<h$. (Take $\rho=\dfrac{\sqrt2}{\pi}\mu$C/m$^3$ , $h=(\sqrt2+1)$ m)
\end{bankq}
\begin{bankq}[id=TOAS06-P1-PHY-S3-Q06, source={TOAS - Test 06 (17 Apr 2022)}, paper=1, qno=6, type=numerical, subtopic={Gauss's law}, answer={0.00}, answer_src=key, difficulty=3, diff_parts={K1 S1 I1 C0 F1}, diff_note={concepts: Gauss's law | steps: divergence of E is zero so the enclosed charge is zero | I: use divergence}]
Electric field in a region is given by $\vec E=3y\hat i+3x\hat j+5\hat k$. Find out total charge enclosed in a cube of side `10cm' and sides along $x$, $y$ \& $z$ axis as shown :-
\bfig{113.1}{TOAS06_p1_phy_s3_q06_1.png}
\end{bankq}
\begin{bankq}[id=TOAS07-P1-PHY-S1-Q02, source={TOAS - Test 07 (01 May 2022)}, paper=1, qno=2, type=single, subtopic={Electric potential}, answer={C}, answer_src=key, difficulty=4, diff_parts={K2 S2 I1 C0 F0}, diff_note={concepts: angular momentum conservation, energy with shell potential | steps: L about centre gives v = u/2 when grazing; energy equation at the surface; solve for u | I: grazing means velocity tangential at the surface}]
A fixed non-conducting spherical shell of radius $R$ has a charge $Q$ spread uniformly on its surface. A particle having mass `$m$' and charge `$q$' is projected towards the shell from a large distance at a speed $u$. Impact parameter is $\dfrac R2$ as shown in figure. The particle passes the shell grazing its surface. The projection speed $u$ of the particle is (Neglect gravity).
\bfig{92.0}{TOAS07_p1_phy_s1_q02_B.png}
\begin{qoptions}
\opt{A}{$\sqrt{\dfrac{4Qq}{3\pi\epsilon_0 mR}}$}
\opt{B}{$\sqrt{\dfrac{Qq}{2\pi\epsilon_0 mR}}$}
\opt{C}{$\sqrt{\dfrac{2Qq}{3\pi\epsilon_0 mR}}$}
\opt{D}{$\sqrt{\dfrac{3Qq}{2\pi\epsilon_0 mR}}$}
\end{qoptions}
\end{bankq}
\begin{bankq}[id=TOAS07-P1-PHY-S2-Q06, source={TOAS - Test 07 (01 May 2022)}, paper=1, qno=6, type=multi, subtopic={Electric potential}, answer={A,C}, answer_src=key, difficulty=3, diff_parts={K1 S2 I0 C0 F1}, diff_note={concepts: motion in the field of a line charge | steps: potential of a line charge is logarithmic; energy conservation gives speed at eta R; angular momentum for the transverse velocity}]
An infinite line charge is perpendicular to the plane of the figure having linear charge density $\lambda$. A particle having charge $+Q$ and mass $m$ is projected in the field of the line charge from point $P$. The point $P$ is at a distance $R$ from the line charge and velocity given to the particle is $v_0$ perpendicular to the radial line at $P$ as shown in figure.
\bfig{106.0}{TOAS07_p1_phy_s2_q06_1.png}
\begin{qoptions}
\opt{A}{Speed of the particle when its distance from the line charge grows to $\eta R$ ($\eta>1$), is $\sqrt{v_0^2+\dfrac{\lambda Q}{\pi\epsilon_0 m}\ln(\eta)}$}
\opt{B}{Speed of the particle when its distance from the line charge grows to $\eta R$ ($\eta>1$), is $\sqrt{v_0^2-\dfrac{\lambda Q}{\pi\epsilon_0 m}\ln(\eta)}$}
\opt{C}{Velocity component of the particle along the radial line (joining the line charge to the particle) at the instant its distance is $\eta R$, is given by $\sqrt{v_0^2\left(1-\dfrac1{\eta^2}\right)+\dfrac{\lambda Q}{\pi\epsilon_0 m}\ln(\eta)}$}
\opt{D}{Velocity component of the particle along the radial line (joining the line charge to the particle) at the instant its distance is $\eta R$, is given by $\sqrt{v_0^2\left(1+\dfrac1{\eta^2}\right)+\dfrac{\lambda Q}{\pi\epsilon_0 m}\ln(\eta)}$}
\end{qoptions}
\end{bankq}
\begin{bankq}[id=TOAS07-P1-PHY-S2-Q07, source={TOAS - Test 07 (01 May 2022)}, paper=1, qno=7, type=multi, subtopic={Dipoles}, answer={A,C,D}, answer_src=key, difficulty=3, diff_parts={K1 S2 I0 C0 F1}, diff_note={concepts: dipole oscillations in a uniform field | steps: restoring torque qE l theta; moment of inertia m l\^{}2/2; omega and period; energy options}]
Figure shows an electric dipole formed by two particles fixed at the ends of a light rod of length $\ell$. The mass of each particle is $m$ and the charges are $-q$ and $+q$. The system is placed in such a way that the dipole axis is parallel to a uniform electric field $E$ that exist in the region. The dipole is slightly rotated about its centre and released. Considering small angular displacement $\theta$, which of the following options are correct.
\bfig{104.8}{TOAS07_p1_phy_s2_q07_1.png}
\begin{qoptions}
\opt{A}{Motion of the dipole is angular SHM with time period $2\pi\sqrt{\dfrac{m\ell}{2qE}}$}
\opt{B}{Motion of the dipole is angular SHM with time period $2\pi\sqrt{\dfrac{m\ell}{qE}}$}
\opt{C}{During SHM no net force acts on dipole.}
\opt{D}{During SHM, when dipole moment vector becomes parallel to electric field, net torque on dipole is zero.}
\end{qoptions}
\end{bankq}
\begin{bankq}[id=TOAS07-P1-PHY-S2-Q08, source={TOAS - Test 07 (01 May 2022)}, paper=1, qno=8, type=multi, subtopic={Electric potential}, answer={A,C}, answer_src=key, difficulty=4, diff_parts={K1 S2 I1 C0 F1}, diff_note={concepts: equipotential of two point charges | steps: set V = 0 for charges -2Q and Q; Apollonius circle; centre and radius | I: ratio of distances is constant}]
Two fixed charges $-2Q$ and $Q$ are located at the points with coordinates $(-3a,0)$ and $(+3a,0)$ respectively in the $xy$ plane.
\begin{qoptions}
\opt{A}{All the points in $x$-$y$ plane where the electric potential due to the two charges is zero lie on a circle of radius $4a$.}
\opt{B}{All the points in $x$-$y$ plane where the electric potential due to the two charges is zero lie on a circle having its centre at $(5a,0)$}
\opt{C}{All the points in $x$-$y$ plane where the electric potential due the two charges is zero lie on an ellipse.}
\opt{D}{All the points in $x$-$y$ plane where the electric potential due the two charges is zero lie on an ellipse having its centre at $(5a,0)$}
\end{qoptions}
\end{bankq}
\begin{bankpassage}[id=TOAS07-P1-PHY-S3-G3-4, source={TOAS - Test 07 (01 May 2022)}, questions={3-4}]
In a uniformly charged fixed nonconducting sphere, a very thin tunnel has been made along the diameter as shown in figure. A charge particle $-q$ having mass $m$ is released from rest at $t=0$ at one end of the tunnel. Total charge $Q$ of sphere is uniformly distributed in its volume. Charge particle reaches at the centre of solid sphere first time at $t=\pi\sqrt{\dfrac{\alpha.\,\pi\epsilon_0 mR^3}{Q.\,q}}$ with speed $v=\sqrt{\dfrac{Qq}{N.\,\pi\epsilon_0 mR}}$.
\bfig{72.3}{TOAS07_p1_phy_s3_para3_4_1.png}
\end{bankpassage}
\begin{bankq}[id=TOAS07-P1-PHY-S3-Q03, source={TOAS - Test 07 (01 May 2022)}, paper=1, qno=3, type=numerical, subtopic={Electric field}, answer={4.00}, answer_src=key, difficulty=3, diff_parts={K1 S2 I0 C0 F1}, diff_note={concepts: electric field problem from the passage | steps: field of the source charges; condition; value of N}, group=TOAS07-P1-PHY-S3-G3-4]
The value of $N$ is\underline{\hspace{1cm}}
\end{bankq}
\begin{bankq}[id=TOAS07-P1-PHY-S3-Q04, source={TOAS - Test 07 (01 May 2022)}, paper=1, qno=4, type=numerical, subtopic={Electric field}, answer={1.00}, answer_src=key, difficulty=2, diff_parts={K1 S1 I0 C0 F1}, diff_note={concepts: electric field problem from the passage | steps: field of the source charges; value of alpha}, group=TOAS07-P1-PHY-S3-G3-4]
The value of $\alpha$ is \underline{\hspace{1.2cm}}
\end{bankq}
\begin{bankq}[id=TOAS07-P1-PHY-S4-Q09, source={TOAS - Test 07 (01 May 2022)}, paper=1, qno=9, type=numerical, subtopic={Electric field}, answer={2}, answer_src=key, difficulty=3, diff_parts={K1 S1 I1 C0 F1}, diff_note={concepts: oscillation of a charge at the centre of a ring | steps: restoring force on the charge for small displacement along the axis; omega squared varies as 1/R\^{}3 | I: scaling of omega with R}]
A ring of radius $R$ that lies in the $yz$ plane carries a positive charge $Q$ uniformly distributed over its circumference. A particle of mass $m$ that carries a negative charge of magnitude $q$ is at the center of the ring. When the charges $Q$ and $q$ are 5 $\mu$C and $-5$ $\mu$C, respectively, and the radius of the ring is 8.0 cm, the mass $m$ oscillates about its equilibrium position with an angular frequency of 16 rad/s. Find the angular frequency (in rad/s) of oscillation of the mass if the radius of the ring is made four times to 32 cm and all other parameters remain unchanged. (Round off to nearest integer).
\end{bankq}
\begin{bankq}[id=TOAS07-P2-PHY-S1-Q01, source={TOAS - Test 07 (01 May 2022)}, paper=2, qno=1, type=multi, subtopic={Dipoles}, answer={A,B}, answer_src=key, difficulty=4, diff_parts={K1 S2 I1 C0 F1}, diff_note={concepts: torque on a charged ring in a field and SHM | steps: charge centre of the half ring; torque about O; moment of inertia; omega and period | I: treat the arc as a point charge at its centroid}]
A semi-circular ring of mass $m$ and radius $R$ with linear charge density $-\lambda$, hinged at its centre with the help of a massless spoke as shown is placed in a uniform electric field as shown in the figure. If the ring is slightly rotated about $O$ in the plane of paper and released. Take $m=8$kg, $\lambda=2$C/m and $E=2$N/c.
\bfig{44.3}{TOAS07_p2_phy_s1_q01_1.png}
\begin{qoptions}
\opt{A}{ring will perform angular SHM}
\opt{B}{the time period of oscillation is $2\pi$ sec.}
\opt{C}{shown position is unstable equilibrium}
\opt{D}{for large angular displacement motion will still be SHM}
\end{qoptions}
\end{bankq}
\begin{bankq}[id=TOAS07-P2-PHY-S1-Q03, source={TOAS - Test 07 (01 May 2022)}, paper=2, qno=3, type=multi, subtopic={Dipoles}, answer={A,B,D}, answer_src=key, difficulty=2, diff_parts={K1 S1 I0 C0 F1}, diff_note={concepts: field and potential of a dipole | steps: plane perpendicular to the dipole axis; equipotential at far distance}]
An electric dipole is placed along $x$-axis with its centre at origin :
\begin{qoptions}
\opt{A}{On a $y$-$z$ plane passing through origin, the electric field is perpendicular at each and every point on that plane at a point far from dipole.}
\opt{B}{$y$-$z$ plane passing through origin is an equipotential surface}
\opt{C}{$y$-$z$ plane passing through origin is not an equipotential surface}
\opt{D}{Flux of electric field through a closed surface enclosing dipole is zero}
\end{qoptions}
\end{bankq}
\begin{bankpassage}[id=TOAS07-P2-PHY-S3-G1-2, source={TOAS - Test 07 (01 May 2022)}, questions={1-2}]
An infinitely long non-conducting thin rod lying on the $z$-axis carries a uniform charge density $\lambda$ $(>0)$ in the region $z>0$ and $-\lambda$ in the region $z<0$.
\end{bankpassage}
\begin{bankq}[id=TOAS07-P2-PHY-S3-Q01, source={TOAS - Test 07 (01 May 2022)}, paper=2, qno=1, type=numerical, subtopic={Electric field}, answer={2.00}, answer_src=key, difficulty=2, diff_parts={K1 S1 I0 C0 F1}, diff_note={concepts: electric field of a line charge distribution | steps: use Gauss's law or integrate; value of n}, group=TOAS07-P2-PHY-S3-G1-2]
The electric field in the $x$-$y$ plane due to the charge distribution at a point which is at a distance $r$ from the origin $O$ is $\dfrac{\lambda}{n\pi\epsilon_0r}$. Find $n$.
\end{bankq}
\begin{bankq}[id=TOAS07-P2-PHY-S3-Q02, source={TOAS - Test 07 (01 May 2022)}, paper=2, qno=2, type=numerical, subtopic={Electric field}, answer={4.00}, answer_src=key, difficulty=2, diff_parts={K1 S1 I0 C0 F1}, diff_note={concepts: flux through a circular region | steps: field component normal to the plane; integrate over the disc}, group=TOAS07-P2-PHY-S3-G1-2]
The flux due to this electric field passing through the circular region : $x^2+y^2=R^2$ in $xy$ plane is $\dfrac{n\lambda R}{4\epsilon_0}$. Find $n$.
\bfig{57.4}{TOAS07_p2_phy_s3_q02_1.png}
\end{bankq}
\begin{bankq}[id=TOAS07-P2-PHY-S4-Q07, source={TOAS - Test 07 (01 May 2022)}, paper=2, qno=7, type=numerical, subtopic={Electric potential}, answer={6}, answer_src=key, difficulty=4, diff_parts={K1 S2 I0 C1 F1}, diff_note={concepts: energy conservation with potential of two charges | steps: potential at C and at D; kinetic energy becomes zero at D; solve for the distance}]
Two fixed, equal, positive charges, each of magnitude $q=5\times10^{-5}$ C are located at points A and B separated by a distance of 6m. An equal and opposite charge moves towards them along the line COD, the perpendicular bisector of the line AB. The moving charge, when it reaches the point C at a distance of 4m from O, has a kinetic energy of 1.5 joules. Calculate the distance of the farthest point D from A (in m) which the negative charge will reach before returning towards C.
\bfig{154.3}{TOAS07_p2_phy_s4_q07_1.png}
\end{bankq}
\begin{bankq}[id=TOAS07-P2-PHY-S4-Q08, source={TOAS - Test 07 (01 May 2022)}, paper=2, qno=8, type=numerical, subtopic={Coulomb's law}, answer={4}, answer_src=key, difficulty=4, diff_parts={K1 S2 I1 C0 F1}, diff_note={concepts: orbiting charges and mechanical energy | steps: speed of each charge from the net force; kinetic energy of two; potential energy of the system | I: both charges interact with each other as well}]
Two bodies of charges $-Q$ and mass $m$ each orbit a body of charge $Q$ (at rest) in a circle of radius $r$. What is the total mechanical energy of the system? If answer is $\dfrac{-3kQ^2}{nr}$, find `$n$'. (Here $k=\dfrac{1}{4\pi\epsilon_0}$) (Neglect self energy)
\bfig{67.3}{TOAS07_p2_phy_s4_q08_1.png}
\end{bankq}
\begin{bankq}[id=TOAS08-P1-PHY-S1-Q01, source={TOAS - Test 08 (15 May 2022)}, paper=1, qno=1, type=multi, subtopic={Coulomb's law}, answer={A,B}, answer_src=key, difficulty=4, diff_parts={K2 S2 I0 C0 F1}, diff_note={concepts: Coulomb repulsion with a spring, centre of mass motion | steps: energy conservation in the centre of mass frame gives the maximum separation; non-SHM nature of the oscillation; centre of mass acceleration is zero}]
Two identical balls each of mass $m$ carries equal charge $Q$ are connected to a massless spring of spring constant $k$ and released from rest. Initially spring was in it natural length $\ell$. Take $k=\dfrac{2Q^2}{3\pi\varepsilon_0\ell^3}$, then select correct option(s)
\bfig{50.4}{TOAS08_p1_phy_s1_q01_1.png}
\begin{qoptions}
\opt{A}{Maximum distance between balls in motion will be $\dfrac{3\ell}{2}$.}
\opt{B}{Both ball will perform oscillation in frame of centre of mass but these oscillation will not be simple harmonic.}
\opt{C}{Acceleration of centre of mass of system will be decreasing with time.}
\opt{D}{Acceleration of centre of mass of system will be increasing with time.}
\end{qoptions}
\end{bankq}
\begin{bankq}[id=TOAS08-P1-PHY-S1-Q06, source={TOAS - Test 08 (15 May 2022)}, paper=1, qno=6, type=multi, subtopic={Conductors}, answer={A,B,C,D}, answer_src=key, difficulty=4, diff_parts={K1 S2 I1 C0 F1}, diff_note={concepts: charges on concentric conducting shells | steps: B is earthed so V\_B = 0; A and C connected so V\_A = V\_C; solve for the charges on each surface | I: charge conservation with potential conditions}]
Three conducting spherical shell $A$,$B$ and $C$ has radius $R$, $2R$ and $3R$ respectively. Conductor $B$ is connected to earth and $A$ and $C$ are connected by conducting wire. If charge $Q$ is given to conductor $C$ then :-
\bfig{106.5}{TOAS08_p1_phy_s1_q06_1.png}
\begin{qoptions}
\opt{A}{Charge on outer surface of conductor $A$ is $\dfrac{2Q}{11}$.}
\opt{B}{Charge on inner surface of conductor $C$ is $\dfrac{6Q}{11}$.}
\opt{C}{Charge on conductor $B$ is $-\dfrac{8Q}{11}$.}
\opt{D}{Charge on outer surface of conductor $C$ is $\dfrac{3Q}{11}$}
\end{qoptions}
\end{bankq}
\begin{bankpassage}[id=TOAS08-P1-PHY-S2-G7-8, source={TOAS - Test 08 (15 May 2022)}, questions={7-8}]
A Non-conducting ring of mass $m$ and radius $R$ is placed on a rough horizontal surface as shown in figure.
\bfig{131.5}{TOAS08_p1_phy_s2_para7_8_1.png}
It has linear charge density $\lambda=\lambda_0\sin\phi$ where $\lambda_0$ is constant and angle $\phi$ is angle of radius vector with line $OA$.
\end{bankpassage}
\begin{bankq}[id=TOAS08-P1-PHY-S2-Q07, source={TOAS - Test 08 (15 May 2022)}, paper=1, qno=7, type=single, subtopic={Electric field}, answer={B}, answer_src=key, difficulty=1, diff_parts={K1 S1 I0 C0 F0}, diff_note={concepts: electric field of a nonuniform ring | steps: element field; integrate the components}, group=TOAS08-P1-PHY-S2-G7-8]
Electric field intensity at centre of ring is :
\begin{qoptions}
\opt{A}{$\dfrac{\lambda_0}{8\pi\varepsilon_0R}(\hat i)$}
\opt{B}{$\dfrac{\lambda_0}{4\varepsilon_0R}(-\hat j)$}
\opt{C}{$\dfrac{\lambda_0}{8\varepsilon_0R}(-\hat j)$}
\opt{D}{$\dfrac{\lambda_0}{4\pi^2\varepsilon_0R}(+\hat j)$}
\end{qoptions}
\end{bankq}
\begin{bankq}[id=TOAS08-P1-PHY-S2-Q08, source={TOAS - Test 08 (15 May 2022)}, paper=1, qno=8, type=single, subtopic={Electric field}, answer={C}, answer_src=key, difficulty=4, diff_parts={K2 S2 I1 C0 F0}, diff_note={concepts: rolling motion, torque of electric force | steps: torque of the electric force; energy of the rolling ring; moment of inertia about the contact; speed after pi/2 | I: use energy for a rolling charged ring}, group=TOAS08-P1-PHY-S2-G7-8]
When external uniform electric field $\vec E=E_0\hat i$ is switch on at $t=0$ then speed of ring after it rotates by angel $\pi/2$ will be (assuming friction is sufficient to prevent slipping]
\begin{qoptions}
\opt{A}{$R\sqrt{\dfrac{3\pi\lambda_0E_0}{m}}$}
\opt{B}{$R\sqrt{\dfrac{\lambda_0E_0}{m}}$}
\opt{C}{$R\sqrt{\dfrac{\pi\lambda_0E_0}{m}}$}
\opt{D}{$\pi R\sqrt{\dfrac{2\lambda_0E_0}{m}}$}
\end{qoptions}
\end{bankq}
\begin{bankq}[id=TOAS08-P1-PHY-S3-Q04, source={TOAS - Test 08 (15 May 2022)}, paper=1, qno=4, type=numerical, subtopic={Electric field}, answer={24.00}, answer_src=key, difficulty=4, diff_parts={K1 S1 I2 C0 F1}, diff_note={concepts: force by a point charge on a charged sheet | steps: force on the sheet equals sigma times the flux of the field through it; flux is a fraction of Q/epsilon0 by symmetry | I: flux fraction from the cube symmetry}]
A square non-conducting sheet of side $a$ has uniform charge density $\sigma$. A point charge $Q$ is placed at perpendicular distance $d=a$ above one of it's edges. Then component of electric force acting on sheet perpendicular to it's plane is $F$. Find $\dfrac{Q\sigma}{F\varepsilon_0}$.
\end{bankq}
\begin{bankq}[id=TOAS08-P1-PHY-S3-Q06, source={TOAS - Test 08 (15 May 2022)}, paper=1, qno=6, type=numerical, subtopic={Dipoles}, answer={1.01}, answer_src=key, difficulty=5, diff_parts={K2 S2 I1 C0 F1}, diff_note={concepts: dipole in the field of a point charge, spring equilibrium | steps: interaction energy or force on a dipole on the axis; equilibrium with the spring; solve for the extension | I: force from the field gradient}]
A small block having electric dipole moment $p$ is placed on a frictionless horizontal table is attached to one end of a spring of force constant $k$ as shown. Now a point charge $-q$ is brought very slowly from a great distance towards block along line coinciding the axis of the spring. Extension in spring when block stick to charge is $\left[\dfrac{apq}{b\pi\varepsilon_0k}\right]^{\frac14}$ where $a$ and $b$ are smallest possible integer than $\dfrac{a-b}{100}$ is.
\bfig{129.3}{TOAS08_p1_phy_s3_q06_1.png}
\end{bankq}
\begin{bankq}[id=TOAS08-P2-PHY-S1-Q01, source={TOAS - Test 08 (15 May 2022)}, paper=2, qno=1, type=multi, subtopic={Conductors}, answer={A,B,D}, answer_src=key, difficulty=4, diff_parts={K1 S2 I1 C0 F1}, diff_note={concepts: conducting shell in a field | steps: inner charge induces charge on the inner surface; outer surface charge in the external field; forces on the 4 C charge and the shell | I: shielding by the shell}]
A conducting shell have total charge 6C is placed inside a uniform EF of 80 N/C as shown. Net force on 4C charge particle placed inside shell is $-80\,\mathrm N\,\hat i$ and conductor is held fixed then select the incorrect options :-
\bfig{149.2}{TOAS08_p2_phy_s1_q01_1.png}
\begin{qoptions}
\opt{A}{Force on 4C due to charge induced on outer surface is $80\,\mathrm N\,\hat i$.}
\opt{B}{Force on 4C due to charge induced on outer surface is $320\,\mathrm N\,\hat i$.}
\opt{C}{Force on shell due to 4C is $400\,\mathrm N\,\hat i$.}
\opt{D}{Force on charge particle due to charge induced on inner surface of shell is $+80\,\mathrm N\,\hat i$}
\end{qoptions}
\end{bankq}
\begin{bankq}[id=TOAS08-P2-PHY-S1-Q04, source={TOAS - Test 08 (15 May 2022)}, paper=2, qno=4, type=multi, subtopic={Electric field}, answer={A,C,D}, answer_src=key, difficulty=4, diff_parts={K1 S2 I1 C0 F1}, diff_note={concepts: field of an oppositely charged hemispherical shell | steps: symmetry of the field on the x-axis; inside field direction; equipotential plane | I: symmetry of the charge distribution}]
Charge density on right half of an insulating spherical shell is positive and charge density on left half is negative but same magnitude as on left half. Choose correct options :-
\bfig{102.4}{TOAS08_p2_phy_s1_q04_1.png}
\begin{qoptions}
\opt{A}{Electric field on $x$-axis between $A$ \& $B$ is in same direction at all points.}
\opt{B}{Electric field between $A$ \& $O$ and $O$ \& $B$ are along opposite directions.}
\opt{C}{$y$-$z$ plane inside sphere is equipotential.}
\opt{D}{Electric field on all points on $x$-axis is in same direction outside the sphere.}
\end{qoptions}
\end{bankq}
\begin{bankq}[id=TOAS08-P2-PHY-S1-Q06, source={TOAS - Test 08 (15 May 2022)}, paper=2, qno=6, type=multi, subtopic={Electric field}, answer={A,D}, answer_src=key, difficulty=5, diff_parts={K2 S2 I1 C0 F1}, diff_note={concepts: charged ring with a gap, torque and rotation | steps: ring with a gap equals a full ring plus opposite charge at the gap; torque about the COM; energy gives maximum angular speed | I: superposition with a negative gap charge}]
An uniform insulating ring and radius $R$ has total charge $q$ distributed uniformly along its circumference except on length $d$ ($d\ll R$). It is placed on smooth surface in uniform electric field parallel to length $d$ which is between $A$ \& $B$. Initially ring is released from rest. Select correct options :-
\bfig{101.4}{TOAS08_p2_phy_s1_q06_1.png}
\begin{qoptions}
\opt{A}{Initial torque about COM is $\dfrac{qEd}{2\pi}$.}
\opt{B}{Initial torque about COM is $\dfrac{qEd}{\pi}$}
\opt{C}{Maximum angular speed in subsequent motion will be $\dfrac1R\sqrt{\dfrac{qdE}{2\pi m}}$.}
\opt{D}{Maximum angular speed in subsequent motion will be $\dfrac1R\sqrt{\dfrac{qdE}{\pi m}}$.}
\end{qoptions}
\end{bankq}
\begin{bankq}[id=TOAS08-P2-PHY-S2-Q07, source={TOAS - Test 08 (15 May 2022)}, paper=2, qno=7, type=single, subtopic={Conductors}, answer={A}, answer_src=key, difficulty=2, diff_parts={K1 S2 I0 C0 F0}, diff_note={concepts: concentric conducting shells with a charge between them | steps: induced charges; potentials at the centre and the shells; effect of connecting the shells}]
Between two concentric conducting neutral shells, a charge $Q=1$C is fixed at distance of 2m from centre $O$. Potential at centre is $V_0$, at inner sphere is $V_1$ and outer sphere is $V_2$ and total charge on outer surfaces of inner \& outer spheres are $Q_1$ and $Q_2$ respectively. After connecting both with conducting wire these charges becomes $Q_1'$ and $Q_2'$. Match List--I with List--II.
\bfig{63.7}{TOAS08_p2_phy_s2_q07_1.png}
\begin{center}\small\renewcommand{\arraystretch}{1.1}\setlength{\extrarowheight}{3.5pt}\setlength{\tabcolsep}{4pt}
\begin{tabular}{|c|c|c|c|}
\hline
\multicolumn{2}{|c|}{List-I} & \multicolumn{2}{c|}{List-II}\\[3pt]
\hline
(P) & $\dfrac{V_1}{V_2}$ & (1) & 3\\[6pt]
(Q) & $\left|\dfrac{Q_1}{Q_2}\right|$ & (2) & 1\\[6pt]
(R) & $\left|\dfrac{Q_1'}{Q_2'}\right|$ & (3) & 2\\[6pt]
(S) & $\dfrac{V_0}{V_2}$ & (4) & $\dfrac13$\\[6pt]
 & & (5) & 0\\[3pt]
\hline
\end{tabular}
\end{center}
\begin{qoptions}
\opt{A}{$P\to3;Q\to5;R\to4;S\to3$}
\opt{B}{$P\to2;Q\to2;R\to5;S\to2$}
\opt{C}{$P\to5;Q\to5;R\to5;S\to2$}
\opt{D}{$P\to1;Q\to5;R\to5;S\to1$}
\end{qoptions}
\end{bankq}
\begin{bankq}[id=TOAS08-P2-PHY-S3-Q02, source={TOAS - Test 08 (15 May 2022)}, paper=2, qno=2, type=numerical, subtopic={Gauss's law}, answer={3.73 to 3.75}, answer_src=key, difficulty=5, diff_parts={K1 S2 I1 C1 F1}, diff_note={concepts: Gauss's law with a finite part of a cylinder | steps: charge of the cylinder inside the sphere; integrate the volume; flux = charge over epsilon0 | I: Gauss's law with a finite portion}]
Consider a uniformly charged very long solid non-conducting cylinder of radius $R$ and having volume charge density $\rho=\epsilon_0$. Find electric flux (in SI unit) through a gaussian spherical surface of radius $2R$ and having one of its diameter along the axis of cylinder and centre on axis of cylinder. Take : $R=\left(\dfrac1\pi\right)^{1/3}$
\end{bankq}
\begin{bankq}[id=TOAS09-P1-PHY-S2-Q05, source={TOAS - Test 09 (30 May 2022)}, paper=1, qno=5, type=multi, subtopic={Gauss's law}, answer={A,B,C,D}, answer_src=key, difficulty=4, diff_parts={K1 S2 I1 C0 F1}, diff_note={concepts: flux through cube faces by symmetry | steps: complete the point charge to the centre of a cube; fraction of flux through each face; fractions for each coordinate case | I: build a cube around the charge}]
Two surface OABC and OCDE lies in the plane of xy and yz as shown in the figure. A charged particle `q' lies in the space at a point P, if:
\bfig{110.4}{TOAS09_p1_phy_s2_q05_1.png}
\begin{qoptions}
\opt{A}{coordinates of `P' is $(a-\Delta r,\ a-\Delta r,\ \Delta r)$ and $a\gg\Delta r$, then flux passing through surface OABC is $\dfrac{7q}{24\varepsilon_0}$.}
\opt{B}{coordinates of `P' is $(a-\Delta r,\ a-\Delta r,\ \Delta r)$ and $a\gg\Delta r$, then flux passing through surface OCDE is $\dfrac{q}{24\varepsilon_0}$.}
\opt{C}{coordinates of `P' is $(a+\Delta r,\ a+\Delta r,\ -\Delta r)$ and $a\gg\Delta r$, then flux passing through surface OABC is $\dfrac{q}{24\varepsilon_0}$.}
\opt{D}{coordinates of `P' is $(a+\Delta r,\ a+\Delta r,\ -\Delta r)$ and $a\gg\Delta r$, then flux passing through surface OCDE is $\dfrac{q}{24\varepsilon_0}$.}
\end{qoptions}
\end{bankq}
\begin{bankq}[id=TOAS09-P1-PHY-S2-Q07, source={TOAS - Test 09 (30 May 2022)}, paper=1, qno=7, type=multi, subtopic={Dipoles}, answer={A,B}, answer_src=key, difficulty=5, diff_parts={K2 S2 I1 C0 F1}, diff_note={concepts: dipole moment of a charged shell, rolling dynamics | steps: dipole moment of the half charged shell; torque in the field; rolling constraint with I = 2mR\^{}2/3; accelerations and friction | I: dipole moment from two half charges}]
A non conducting uniform thin shell of mass $m$ and radius $R$ is having uniform charge density $+\sigma$ on one half and $-\sigma$ on another half as shown in figure. It is placed on a rough non conducting horizontal plane. At $t=0$ a uniform electric field $\vec E=-E_0\hat j$ N/C is switched on and the solid sphere starts rolling without slipping. Then which of the following is/are correct.
\bfig{156.7}{TOAS09_p1_phy_s2_q07_1.png}
\begin{qoptions}
\opt{A}{The acceleration of the topmost point of the hollow sphere at $t=0$, is $\dfrac{36}{25}\dfrac{\pi\sigma E_0R^2}{m}$.}
\opt{B}{The magnitude of frictional force at $t=0$ is $\dfrac{18}{25}\pi\sigma E_0R^2$.}
\opt{C}{The acceleration of centre of mass of the hollow sphere at $t=0$, is $\dfrac{12}{25}\dfrac{\pi\sigma E_0R^2}{m}$.}
\opt{D}{The magnitude of frictional force at $t=0$ is $\dfrac{9}{25}\pi\sigma E_0R^2$.}
\end{qoptions}
\end{bankq}
\begin{bankq}[id=TOAS09-P1-PHY-S4-Q08, source={TOAS - Test 09 (30 May 2022)}, paper=1, qno=8, type=numerical, subtopic={Gauss's law}, answer={3}, answer_src=key, difficulty=2, diff_parts={K1 S1 I0 C0 F1}, diff_note={concepts: field inside a non-uniform sphere | steps: E(r) from Gauss's law is proportional to r\^{}(n+1); ratio of fields gives n}]
A non-conducing solid sphere has a charge density given by $\rho=\rho_0r^n$, where $\rho_0$, $n$ are constant and `r' is the distance from the centre of sphere. If the magnitude of electric field at half of the radius is $\dfrac{1}{16}$ times that of the surfaces of nonconducting then find the value of $n$.
\end{bankq}
\begin{bankq}[id=TOAS09-P2-PHY-S1-Q09, source={TOAS - Test 09 (30 May 2022)}, paper=2, qno=9, type=single, subtopic={Gauss's law}, answer={B}, answer_src=key, difficulty=1, diff_parts={K1 S1 I0 C0 F0}, diff_note={concepts: flux through part of a sphere | steps: fraction of the solid angle; flux}]
Find out flux passing through shaded surface of the sphere when a point charge `q' is placed at centre of sphere. [Radius of sphere is $R$]
\bfig{59.6}{TOAS09_p2_phy_s1_q09_1.png}
\begin{qoptions}
\opt{A}{$\dfrac{q}{\varepsilon_0}$}
\opt{B}{$\dfrac{q}{4\varepsilon_0}$}
\opt{C}{$\dfrac{q}{2\varepsilon_0}$}
\opt{D}{Zero.}
\end{qoptions}
\end{bankq}
\begin{bankq}[id=TOAS09-P2-PHY-S1-Q10, source={TOAS - Test 09 (30 May 2022)}, paper=2, qno=10, type=single, subtopic={Gauss's law}, answer={B}, answer_src=key, difficulty=1, diff_parts={K1 S1 I0 C0 F0}, diff_note={concepts: flux through a cube | steps: E\_x on the two faces at x = l and x = 2l (given); flux difference}]
A cube of side `$\ell$' placed as shown in figure $E_x=\beta x^{\frac12}$. Flux passing through cube will be:-
\bfig{112.0}{TOAS09_p2_phy_s1_q10_1.png}
\begin{qoptions}
\opt{A}{$2\sqrt2\,\beta\ell^{\frac52}$}
\opt{B}{$(\sqrt2-1)\,\beta\ell^{\frac52}$}
\opt{C}{$-\beta\ell^{\frac52}$}
\opt{D}{Zero}
\end{qoptions}
\end{bankq}
\begin{bankq}[id=TOAS09-P2-PHY-S1-Q11, source={TOAS - Test 09 (30 May 2022)}, paper=2, qno=11, type=single, subtopic={Electric potential}, answer={A}, answer_src=key, difficulty=3, diff_parts={K1 S2 I1 C0 F0}, diff_note={concepts: potential difference between two charged spheres | steps: potential at the centre of one sphere from itself; from the other sphere; subtract | I: potential of a uniform sphere at its centre}]
Two solid spheres having volume charge density $\rho$ and $-\rho$ respectively as shown in figure. Find out the potential difference $(\Delta V)$ between centre of both sphere for $d=R$
\bfig{101.1}{TOAS09_p2_phy_s1_q11_1.png}
\begin{qoptions}
\opt{A}{$\dfrac{\rho}{3\varepsilon_0}R^2$}
\opt{B}{$\dfrac{\rho}{\varepsilon_0}R^2$}
\opt{C}{0}
\opt{D}{$\dfrac{2\rho}{\varepsilon_0}R^2$}
\end{qoptions}
\end{bankq}
\begin{bankq}[id=TOAS09-P2-PHY-S1-Q12, source={TOAS - Test 09 (30 May 2022)}, paper=2, qno=12, type=single, subtopic={Electric potential}, answer={B}, answer_src=key, difficulty=2, diff_parts={K1 S1 I1 C0 F0}, diff_note={concepts: interaction energy in an infinite chain | steps: series of 1/n with alternating signs; sum is ln 2 | I: alternating harmonic series}]
There is a infinite chain of charge $q$ \& $-q$. Distance between the two neighbouring charge is equal to `d'. The interaction energy of one of charge with all other charge is :
\bfig{150.6}{TOAS09_p2_phy_s1_q12_1.png}
\begin{qoptions}
\opt{A}{$\dfrac{-2q^2}{4\pi\varepsilon_0d}$}
\opt{B}{$\dfrac{-2q^2\ell n2}{4\pi\varepsilon_0d}$}
\opt{C}{$\dfrac{2q^2\ell n2}{4\pi\varepsilon_0d}$}
\opt{D}{$\dfrac{q^2}{\pi\varepsilon_0d}$}
\end{qoptions}
\end{bankq}
\begin{bankq}[id=TOAS09-P2-PHY-S1-Q13, source={TOAS - Test 09 (30 May 2022)}, paper=2, qno=13, type=single, subtopic={Electric potential}, answer={A}, answer_src=key, difficulty=1, diff_parts={K1 S1 I0 C0 F0}, diff_note={concepts: potential of arcs at a common point | steps: each arc is at distance R from C; V = k times total charge over R}]
Three arc are joined together as shown in figure having charges $Q$, $-2Q$, $3Q$ respectively. Find out the electric potential at point C.
\bfig{100.9}{TOAS09_p2_phy_s1_q13_1.png}
\begin{qoptions}
\opt{A}{$\dfrac{Q}{2\pi\varepsilon_0R}$}
\opt{B}{$\dfrac{5Q}{12\pi\varepsilon_0R}$}
\opt{C}{$\dfrac{3Q}{32\pi\varepsilon_0R}$}
\opt{D}{$\dfrac{3Q}{8\pi\varepsilon_0R}$}
\end{qoptions}
\end{bankq}
\begin{bankq}[id=TOAS09-P2-PHY-S1-Q18, source={TOAS - Test 09 (30 May 2022)}, paper=2, qno=18, type=single, subtopic={Conductors}, answer={C}, answer_src=key, difficulty=2, diff_parts={K1 S1 I1 C0 F0}, diff_note={concepts: potential outside a conductor with a cavity charge | steps: shell induces charges; outside the sphere the field is that of the charge at the centre | I: shielding by the conductor}]
A point charge $-q$ is placed inside cavity of a metallic sphere of radius $R$. Radius of the cavity is $r$. Potential at point P will be -
\bfig{139.1}{TOAS09_p2_phy_s1_q18_1.png}
\begin{qoptions}
\opt{A}{$\dfrac{-q}{4\pi\varepsilon_0y}$}
\opt{B}{$\dfrac{-q}{4\pi\varepsilon_0y}+\dfrac{q}{4\pi\varepsilon_0x}$}
\opt{C}{$\dfrac{-q}{4\pi\varepsilon_0x}$}
\opt{D}{$\dfrac{q}{4\pi\varepsilon_0y}-\dfrac{q}{4\pi\varepsilon_0x}$}
\end{qoptions}
\end{bankq}
\begin{bankq}[id=TOAS09-P2-PHY-S2-Q03, source={TOAS - Test 09 (30 May 2022)}, paper=2, qno=3, type=numerical, subtopic={Electric potential}, answer={0}, answer_src=solved, working={$V=\frac1{4\pi\varepsilon_0}\int_0^{2\pi}\lambda_0\cos\frac\theta2\,d\theta=\frac{\lambda_0}{4\pi\varepsilon_0}\left[2\sin\frac\theta2\right]_0^{2\pi}=0$}, difficulty=2, diff_parts={K1 S1 I0 C0 F1}, diff_note={concepts: potential of a nonuniform ring | steps: integral of lambda0 cos(theta/2) over the ring is zero}]
Linear charge density of ring of radius `R' given as $\lambda=\lambda_0\cos\dfrac\theta2$. The potential at centre of ring will be
\bfig{99.7}{TOAS09_p2_phy_s2_q03_1.png}
\end{bankq}
\begin{bankq}[id=TOAS09-P2-PHY-S2-Q04, source={TOAS - Test 09 (30 May 2022)}, paper=2, qno=4, type=numerical, subtopic={Electric potential}, answer={100}, answer_src=solved, working={$V_A-V_O=-\vec E\cdot\vec{OA}$ with $\vec{OA}=-2\hat i$ and $E_x=100\cos60^\circ$... equals $+100$ V}, difficulty=2, diff_parts={K1 S1 I0 C0 F1}, diff_note={concepts: potential difference in a uniform field | steps: V\_A - V\_O = -E d cos(angle)}]
Uniform $\vec E=100$v/m is directed as shown in figure. The potential difference between $v_A-v_O$ is
\bfig{84.9}{TOAS09_p2_phy_s2_q04_1.png}
\end{bankq}
\begin{bankq}[id=TOAS09-P2-PHY-S2-Q10, source={TOAS - Test 09 (30 May 2022)}, paper=2, qno=10, type=numerical, subtopic={Conductors}, answer={2}, answer_src=key, difficulty=4, diff_parts={K1 S2 I1 C0 F1}, diff_note={concepts: earthing of concentric conductors | steps: induced charges on the thick shell; potential of the outer shell is zero; charge flowing to earth | I: potential of the earthed shell is zero}]
A thick conducting shell of inner radius $R$ and outer radius $2R$ has charge $Q$ at its center. It is covered by a thin conducting shell of radius $4R$ which has charge $2Q$. If charge flown into the earth when switch is closed is $nQ$, find $n$.
\bfig{106.5}{TOAS09_p2_phy_s2_q10_1.png}
\end{bankq}
\begin{bankq}[id=TOAS10-P1-PHY-S1-Q06, source={TOAS - Test 10 (12 Jun 2022)}, paper=1, qno=6, type=multi, subtopic={Conductors}, answer={A,C}, answer_src=key, difficulty=3, diff_parts={K1 S1 I1 C0 F1}, diff_note={concepts: conducting sphere and the method of images | steps: earthed sphere acts as image charge -qR/r at R\^{}2/r; force from the image charge | I: image charge}]
A charge q is placed at a distance r from centre of conducting sphere as shown in figure. Now sphere is earthed then correct options are :-
\bfig{109.7}{TOAS10_p1_phy_s1_q06_1.png}
\begin{qoptions}
\opt{A}{Induced charge on sphere is $-\left(\dfrac Rr\right)q$}
\opt{B}{Induced charge on sphere is $-q$}
\opt{C}{Force acting between sphere and external charge q is $\dfrac{1}{4\pi\varepsilon_0}\dfrac{q^2Rr}{(r^2-R^2)^2}$}
\opt{D}{Force between sphere and external charge q is $\dfrac{q^2}{r^2}$.}
\end{qoptions}
\end{bankq}
\begin{bankq}[id=TOAS10-P1-PHY-S3-Q01, source={TOAS - Test 10 (12 Jun 2022)}, paper=1, qno=1, type=numerical, subtopic={Electric field}, answer={9.00}, answer_src=key, difficulty=4, diff_parts={K1 S2 I1 C0 F1}, diff_note={concepts: projectile with an additional constant force | steps: height from the uncharged case; effective acceleration with +q; acceleration with -q and the speed | I: effective gravity with the field force}]
A particle is uncharged and is thrown vertically upward from ground level with a speed of $5\sqrt5$ m/s. As a result, it attains a maximum height $h$. The particle is then given a positive charge $+q$ and reaches the same maximum height $h$ when thrown vertically upward with a speed of 13 m/s. Finally, the particle is given a negative charge $-q$. Ignoring air resistance, determine the speed (in m/s) with which the negatively charged particle must be thrown vertically upward, so that it attains exactly the same maximum height $h$.
\end{bankq}
\begin{bankq}[id=TOAS10-P1-PHY-S3-Q02, source={TOAS - Test 10 (12 Jun 2022)}, paper=1, qno=2, type=numerical, subtopic={Electric potential}, answer={2.00}, answer_src=key, difficulty=3, diff_parts={K1 S1 I1 C0 F1}, diff_note={concepts: work by a nonuniform electric field on a ring | steps: E dot dr = alpha d(xy) so the work is q alpha times the change of xy; speed from kinetic energy | I: the field is conservative}]
A smooth parabolic wire track defined by relation $y=\dfrac{x^2}{2a}$ is fixed in x-y plane in a gravity free space as shown in figure. An electric field $\vec E=\alpha y\hat i+\alpha x\hat j$ (where $\alpha$ is a constant) exist in the region. A small ring the mass `m', charge `q', which just fits in wire track, is released from rest at position $A\left(-a,\dfrac a2\right)$. Find speed of the ring when it reaches at point $B\left(a,\dfrac a2\right)$ if $q\alpha=\dfrac{2m}{a^2}$.
\bfig{77.3}{TOAS10_p1_phy_s3_q02_1.png}
\end{bankq}
\begin{bankq}[id=TOAS10-P1-PHY-S3-Q03, source={TOAS - Test 10 (12 Jun 2022)}, paper=1, qno=3, type=numerical, subtopic={Electric potential}, answer={0.38 to 0.39}, answer_src=key, difficulty=4, diff_parts={K1 S2 I0 C1 F1}, diff_note={concepts: electrostatic energy of concentric shells | steps: potential of each shell from the charges; energy as half the sum of charge times potential; evaluate n}]
Consider four concentric conducting spherical shells of radius $a$, $2a$, $3a$ and $4a$ having charges $Q$, $-Q$, $2Q$ and $Q$ respectively as shown in figure. Total electrostatic potential energy of this charge configuration is $\left(\dfrac{nQ^2}{\pi\varepsilon_0a}\right)$. Find $n$.
\bfig{98.5}{TOAS10_p1_phy_s3_q03_1.png}
\end{bankq}
\begin{bankq}[id=TOAS11-P1-PHY-S1-Q01, source={TOAS - Test 11 (27 Jun 2022)}, paper=1, qno=1, type=single, subtopic={Conductors}, answer={A}, answer_src=key, difficulty=2, diff_parts={K1 S1 I1 C0 F0}, diff_note={concepts: conducting shell with an external point charge | steps: induced charges -q on the inner surface and +q on the outer surface; potential at the centre is the sum of the three contributions | I: potential of induced surface charges at the centre}]
A point charge $q$ is at a distance $r$ from the centre O of an uncharged spherical conducting layer, whose inner and outer radii are equal to $a$ and $b$ respectively. If $r<a$, then the potential at the centre O is:
\bfig{63.4}{TOAS11_p1_phy_s1_q01_1.png}
\begin{qoptions}
\opt{A}{$\dfrac{q}{4\pi\varepsilon_0}\left(\dfrac1r-\dfrac1a+\dfrac1b\right)$}
\opt{B}{$\dfrac{q}{4\pi\varepsilon_0 r}$}
\opt{C}{$\dfrac{q}{4\pi\varepsilon_0}\left(\dfrac1r+\dfrac1a-\dfrac1b\right)$}
\opt{D}{$\dfrac{q}{4\pi\varepsilon_0}\left(\dfrac1a+\dfrac1b-\dfrac1r\right)$}
\end{qoptions}
\end{bankq}
\begin{bankq}[id=TOAS11-P1-PHY-S2-Q05, source={TOAS - Test 11 (27 Jun 2022)}, paper=1, qno=5, type=multi, subtopic={Electric potential}, answer={A,C}, answer_src=key, difficulty=5, diff_parts={K1 S2 I1 C1 F1}, diff_note={concepts: potential and energy of a charged disc | steps: potential at the rim by direct integration; self energy from half of sigma V integrated over the disc; evaluate | I: choose polar coordinates about the rim point}]
A thin non-conducting disc of radius $R$ is uniformly charged with surface charge density $\sigma$. If potential at any point on the circumference of disc be $V$ and total electrostatic potential energy of the disc be $U$, then:
\begin{qoptions}
\opt{A}{$V=\dfrac{\sigma R}{\pi\varepsilon_0}$}
\opt{B}{$V=\dfrac{2\sigma R}{\pi\varepsilon_0}$}
\opt{C}{$U=\dfrac{2\sigma^2R^3}{3\varepsilon_0}$}
\opt{D}{$U=\dfrac43\dfrac{\sigma^2R^3}{\varepsilon_0}$}
\end{qoptions}
\end{bankq}
\begin{bankq}[id=TOAS11-P1-PHY-S2-Q09, source={TOAS - Test 11 (27 Jun 2022)}, paper=1, qno=9, type=multi, subtopic={Conductors}, answer={A,B,C}, answer_src=key, difficulty=4, diff_parts={K1 S2 I1 C0 F1}, diff_note={concepts: charge distribution on parallel plates | steps: outer surfaces carry half of the total; inner surfaces are equal and opposite to neighbours; assign each surface | I: outer surfaces equal}]
Three very large plates are given charges as shown in the figure. If the cross-sectional area of each plate is same, then the final charge distribution on the surface of the plates a, b, c, d, e, f is:
\bfig{92.8}{TOAS11_p1_phy_s2_q09_1.png}
\begin{qoptions}
\opt{A}{7Q on surface e and 3Q on surface f}
\opt{B}{$-$2Q on surface b and 3Q on surface a}
\opt{C}{$-$7Q on surface d and 2Q on surface c}
\opt{D}{The magnitude of charges at surface b, c, d, e is equal}
\end{qoptions}
\end{bankq}
\begin{bankq}[id=TOAS11-P2-PHY-S1-Q08, source={TOAS - Test 11 (27 Jun 2022)}, paper=2, qno=8, type=multi, subtopic={Coulomb's law}, answer={B,C}, answer_src=key, difficulty=5, diff_parts={K2 S2 I1 C0 F1}, diff_note={concepts: charged balls in a groove, elastic collision | steps: electrostatic energy and energy conservation for the speeds; momentum and energy through the collision; maximum speed after the first collision | I: balls start at opposite points and meet by repulsion or attraction}]
On a horizontal table, there is a smooth circular groove of mean radius $R$. The walls of the groove are non conducting. Two metal balls A and B (each having mass $m$ and radius $r$) are placed inside the groove with their centers $R$ apart. The balls just fit inside the groove. The two balls are given charge $+3q$ and $-q$ are released from state of rest as shown in figure. Ignore the non-uniformity in charge distribution as the balls come close together and collide. The collision is elastic.
\bfig{118.7}{TOAS11_p2_phy_s1_q08_1.png}
\begin{qoptions}
\opt{A}{Speed of ball A is just before first collision is $\sqrt{\dfrac{3q^2}{8\pi\epsilon_0m}\left(\dfrac1{2r}-\dfrac1R\right)}$}
\opt{B}{Speed of ball B is just before first collision is $\sqrt{\dfrac{3q^2}{4\pi\epsilon_0m}\left(\dfrac1{2r}-\dfrac1R\right)}$}
\opt{C}{Maximum speed acquired by each ball after they collide for the first time is $\sqrt{\dfrac{q^2}{8\pi\epsilon_0mrR}(4R-7r)}$}
\opt{D}{Maximum speed acquired by each ball after they collide for the first time is $\sqrt{\dfrac{q^2}{4\pi\epsilon_0mrR}(4R-7r)}$}
\end{qoptions}
\end{bankq}
\begin{bankq}[id=TOAS11-P2-PHY-S2-Q09, source={TOAS - Test 11 (27 Jun 2022)}, paper=2, qno=9, type=single, subtopic={Conductors}, answer={A}, answer_src=key, difficulty=2, diff_parts={K1 S2 I0 C0 F0}, diff_note={concepts: cavity in a conductor and charges placed inside and outside | steps: direction of the field in the cavity; potential at A; potential at C for each configuration}]
\textbf{Answer the following by appropriately matching the lists based on the information given in the paragraph.}\par
There is a spherical cavity of radius $R/3$ in a spherical solid conductor of radius $R$. A is centre of solid conductor, B is centre of cavity and C is a point outside at a distance $R/3$ from the surface. In List-I, there are certain charge configuration suggested and List-II discusses the electric field and potential.
\bfig{111.7}{TOAS11_p2_phy_s2_q09_1.png}
\smallskip
\noindent{\footnotesize\setlength{\tabcolsep}{2pt}\begin{tabular}{|L{.95cm}|L{2.7cm}|L{.95cm}|L{3.1cm}|}\hline
\multicolumn{2}{|c|}{List-I} & \multicolumn{2}{c|}{List-II}\\\hline
(I) & Charge $+Q$ at B & (P) & Electric field in cavity near A is to left\\\hline
(II) & Charge $+Q$ at C & (Q) & Electric field in cavity near A is to right\\\hline
(III) & Charge $-Q$ at B & (R) & Electric potential at A is $+$ve\\\hline
(IV) & Charge $-Q$ at C and $+Q$ at B & (S) & Electric potential at A is $-$ve\\\hline
 & & (T) & Electric potential at C is $+$ve\\\hline
 & & (U) & Electric potential at C is $-$ve\\\hline
\end{tabular}}
\begin{qoptions}
\opt{A}{I $\to$ P,R,T;II $\to$ R,T;III $\to$ Q,S,U;IV $\to$ P,R,U}
\opt{B}{I $\to$ P,R,T;II $\to$ R,T;III $\to$ Q,S,U;IV $\to$ Q,R,U}
\opt{C}{I $\to$ P,S,T;II $\to$ S,T;III $\to$ Q,R,U;IV $\to$ P,S,U}
\opt{D}{I $\to$ P,R,S;II $\to$ T;III $\to$ Q,S,U;IV $\to$ Q,R,U}
\end{qoptions}
\end{bankq}
\begin{bankq}[id=TOAS11-P2-PHY-S2-Q10, source={TOAS - Test 11 (27 Jun 2022)}, paper=2, qno=10, type=single, subtopic={Conductors}, answer={D}, answer_src=key, difficulty=1, diff_parts={K1 S1 I0 C0 F0}, diff_note={concepts: cavity in a conductor and changing the configuration | steps: effect of each change on the field and potential at C and in the cavity; induced charges}]
\textbf{Answer the following by appropriately matching the lists based on the information given in the paragraph.}\par
There is a spherical cavity of radius $R/3$ in a spherical solid conductor of radius $R$. A is centre of solid conductor, B is centre of cavity and C is a point outside at a distance $R/3$ from the surface. In List-I, there are certain charge configuration suggested and List-II discusses the electric field and potential. Now a charge $+Q$ is placed at B.
\bfig{111.4}{TOAS11_p2_phy_s2_q10_1.png}
\smallskip
\noindent{\footnotesize\setlength{\tabcolsep}{2pt}\begin{tabular}{|L{.95cm}|L{2.7cm}|L{.95cm}|L{3.1cm}|}\hline
\multicolumn{2}{|c|}{List-I} & \multicolumn{2}{c|}{List-II}\\\hline
(I) & Charge is moved slightly to left & (P) & Electric field at C changes.\\\hline
(II) & Charge at B is made $+2Q$ & (Q) & Electric potential at C increases.\\\hline
(III) & Charge at B is made $-Q$ & (R) & Electric potential at C decreases.\\\hline
(IV) & An external uniform electric field is switched on in the direction C to A. & (S) & Electric field at A inside cavity changes.\\\hline
 & & (T) & Electric potential at B due to charges induced on conductor changes\\\hline
 & & (U) & Electric field at B due to charges induced on conductor changes.\\\hline
\end{tabular}}
\begin{qoptions}
\opt{A}{I $\to$ P,T;II $\to$ P,Q,S,T;III $\to$ P,R,S,U;IV $\to$ P,R,S,T}
\opt{B}{I $\to$ Q,U;II $\to$ P,Q,S,U;III $\to$ P,R,T,U;IV $\to$ P,R,S,T,U}
\opt{C}{I $\to$ S,T;II $\to$ P,Q,S,T,U;III $\to$ P,R,S,T,U;IV $\to$ P,Q,T,U}
\opt{D}{I $\to$ S,U;II $\to$ P,Q,S,T;III $\to$ P,R,S,T;IV $\to$ P,R,T,U}
\end{qoptions}
\end{bankq}
\begin{bankq}[id=TOAS12-P1-PHY-S1-Q01, source={TOAS - Test 12 (10 Jul 2022)}, paper=1, qno=1, type=multi, subtopic={Gauss's law}, answer={A,C,D}, answer_src=key, difficulty=3, diff_parts={K1 S2 I0 C0 F1}, diff_note={concepts: flux and field of a point charge at the centre of hemispherical shells | steps: flux through each hemisphere is half of q0/epsilon0; fields at the two radii by Coulomb's law; ratio relations; compare options}]
A charge particle $q_0$ is kept at centre of 2 non conducting hemi spherical shell of radius $R_1$ \& $R_2$. If $E_1$ \& $E_2$ is electric field at surfaces of A \& B then
\bfig{77.9}{TOAS12_p1_phy_s1_q01_1.png}
\begin{qoptions}
\opt{A}{Ratio of flux $\dfrac{\phi_1}{\phi_2}=1$.}
\opt{B}{$E_1=E_2$}
\opt{C}{$\dfrac{R_1}{R_2}=\sqrt{\dfrac{E_2}{E_1}}$}
\opt{D}{$E_1\ne E_2$}
\end{qoptions}
\end{bankq}
\begin{bankq}[id=TOAS12-P1-PHY-S1-Q08, source={TOAS - Test 12 (10 Jul 2022)}, paper=1, qno=8, type=multi, subtopic={Conductors}, answer={A,C}, answer_src=key, difficulty=5, diff_parts={K1 S2 I1 C1 F1}, diff_note={concepts: induced charges on a conducting sphere by point charges | steps: image of each charge in the sphere; fields of the induced charges at P; potential of the images | I: method of images for a sphere}]
An uncharged solid spherical conductor of radius $R$ has its centre at the origin of a cartesian frame OXY Z. Two point charges each of $+Q$ are placed at fixed positions at ($5R/2$, 0, 0) and (0, 0, $3R/2$). There will be induced charges on the sphere due to these charges. Choose the correct option(s)
\bfig{112.5}{TOAS12_p1_phy_s1_q08_1.png}
\begin{qoptions}
\opt{A}{Electric field at $\left(\dfrac R2,\dfrac R2,\dfrac R2\right)$ due to induced charges is given by\par $E_p=-\dfrac{2^{3/2}Q}{(4\pi\epsilon_0R^2)}$ $\left[\left(\dfrac92-2\sqrt3\right)\hat i\right.$ $\left.+\left(\dfrac92+\sqrt3\right)\hat j\right.$ $\left.+(9-\sqrt3)\hat k\right]$}
\opt{B}{Electric field at $\left(\dfrac R2,\dfrac R2,\dfrac R2\right)$ due to induced charges is given by\par $E_p=-\dfrac{2^{3/2}Q}{(4\pi\epsilon_0R^2)}$ $\left[\left(\dfrac92-2\sqrt3\right)\hat i\right.$ $\left.+\left(\dfrac92+\sqrt3\right)\hat j\right.$ $\left.+(9+\sqrt3)\hat k\right]$}
\opt{C}{The potential at P due to induced charges is given by\par $V=\dfrac{Q}{(4\pi\epsilon_0R)\left(\frac{16}{15}-\left(\frac29\right)^{1/2}-\left(\frac23\right)^{1/2}\right)}$}
\opt{D}{The potential at P due to induced charges is given by\par $V=\dfrac{Q}{(4\pi\epsilon_0R)\left(\frac{16}{15}+\left(\frac29\right)^{1/2}+\left(\frac23\right)^{1/2}\right)}$}
\end{qoptions}
\end{bankq}
\begin{bankq}[id=TOAS12-P1-PHY-S1-Q09, source={TOAS - Test 12 (10 Jul 2022)}, paper=1, qno=9, type=multi, subtopic={Gauss's law}, answer={A,B,C,D}, answer_src=key, difficulty=3, diff_parts={K1 S2 I0 C0 F1}, diff_note={concepts: flux through a moving or growing sphere | steps: flux versus time for uniform line charge and a moving sphere; for charge proportional to x; for a growing sphere}]
The graph shows variation of flux ($\phi$) with time ($t$) passing through a sphere for varries conditions. Choose the correct option(s) : (Assume : linear charge on x-axis)
\bfig{98.4}{TOAS12_p1_phy_s1_q09_1.png}
\begin{qoptions}
\opt{A}{Graph `a' may corresponds to the case when a sphere is moving along x-axis with constant speed and the charge distribution is uniform in x-direction.}
\opt{B}{Graph `b' may corresponds to the case when a sphere is moving along x-axis with constant speed and the charge distribution is proportional to $x$.}
\opt{C}{Graph `c' may corresponds to the case when a sphere is moving along x-axis with constant speed and the charge distribution is varying linearly with $x$.}
\opt{D}{Graph `d' may corresponds to the case when a sphere is located at origin and its radius is increasing with uniform rate and charge distribution is proportional to $x$.}
\end{qoptions}
\end{bankq}
\begin{bankq}[id=TOAS12-P1-PHY-S1-Q10, source={TOAS - Test 12 (10 Jul 2022)}, paper=1, qno=10, type=multi, subtopic={Electric potential}, answer={A}, answer_src=key, difficulty=4, diff_parts={K1 S2 I1 C0 F1}, diff_note={concepts: interaction energy of a ring and a rod | steps: potential on the axis of the ring; integrate over the rod; evaluate the logarithm | I: ring potential on its axis}]
A non conducting rod AB of length $\sqrt3R$, uniformly distributed charge of linear charge density $\lambda$ and a non-conducting ring of uniformly distributed charge $Q$, are placed as shown in the figure. Point A is the centre of ring and line AB is the axis of ring, perpendicular to plane of ring. The electrostatic interaction energy between ring and rod is:
\bfig{105.8}{TOAS12_p1_phy_s1_q10_1.png}
\begin{qoptions}
\opt{A}{$\dfrac{Q\lambda}{4\pi\varepsilon_0}\ln(2+\sqrt3)$}
\opt{B}{$\dfrac{Q\lambda}{2\pi\varepsilon_0}\ln(2+\sqrt3)$}
\opt{C}{$\dfrac{Q\lambda}{4\pi\varepsilon_0}\ln(2-\sqrt3)$}
\opt{D}{$\dfrac{Q\lambda}{2\pi\varepsilon_0}\ln(2-\sqrt3)$}
\end{qoptions}
\end{bankq}
\begin{bankq}[id=TOAS12-P1-PHY-S2-Q06, source={TOAS - Test 12 (10 Jul 2022)}, paper=1, qno=6, type=numerical, subtopic={Electric field}, answer={13.95 to 13.96}, answer_src=key, difficulty=4, diff_parts={K1 S2 I1 C0 F1}, diff_note={concepts: oscillation of a point charge on the axis of a ring | steps: field of a ring near its centre on the axis; two-body oscillation with reduced mass; period | I: ring and charge both move}]
Figure shows a ring of mass `M' and radius `R' uniformly charged with charge $Q$. A point charge `$-q$' of mass $m$ is kept at the centre of ring. Ring is free to move. Find the time period of small oscillation of the point charge. [Take : $m=1$ kg, $M=2$kg, $q=-3\mu$C, $Q=40\ \mu$C, $R=2$m]
\bfig{67.5}{TOAS12_p1_phy_s2_q06_1.png}
\end{bankq}
\begin{bankq}[id=TOAS12-P2-PHY-S1-Q03, source={TOAS - Test 12 (10 Jul 2022)}, paper=2, qno=3, type=single, subtopic={Electric potential}, answer={A}, answer_src=key, difficulty=3, diff_parts={K1 S2 I1 C0 F0}, diff_note={concepts: energy and momentum in Coulomb repulsion | steps: energy conservation; at the closest approach the speed is minimal in the sense of radial motion; velocity condition fixes the separation | I: use angular momentum when the velocity is tangential}]
In the diagram shown, the charge $+Q$ is fixed. Another charge $+2q$, is projected from a distance $R$ from the fixed charge. The minimum separation between the two charges when the velocity becomes $\dfrac{1}{\sqrt3}$ times of the projected velocity at that instant, is (Assume gravity to the absent)
\bfig{90.6}{TOAS12_p2_phy_s1_q03_1.png}
\begin{qoptions}
\opt{A}{$\dfrac{\sqrt3}{2}R$}
\opt{B}{$\sqrt3R$}
\opt{C}{$\dfrac12R$}
\opt{D}{$4R$}
\end{qoptions}
\end{bankq}
\begin{bankq}[id=TOAS12-P2-PHY-S3-Q05, source={TOAS - Test 12 (10 Jul 2022)}, paper=2, qno=5, type=numerical, subtopic={Electric potential}, answer={300.00}, answer_src=key, difficulty=5, diff_parts={K2 S2 I1 C0 F1}, diff_note={concepts: potential of a sphere, superposition | steps: potential of the full sphere; remove the cavity as a negative sphere; potential at the sphere centre | I: superposition of a negative cavity}]
A solid sphere of radius `R' is uniformly charged with charge density $\rho$ in its volume. A spherical cavity of radius $\dfrac R2$ is made in the sphere as shown in the figure. Electric potential at the centre of the sphere is $\dfrac{A\rho R^2}{B\epsilon_0}$. Where A and B are positive integers. Find $A^2B$.
\bfig{58.3}{TOAS12_p2_phy_s3_q05_1.png}
\end{bankq}
\begin{bankq}[id=TOAS13-P1-PHY-S1-Q03, source={TOAS - Test 13 (24 Jul 2022)}, paper=1, qno=3, type=multi, subtopic={Gauss's law}, answer={B,C}, answer_src=key, difficulty=6, diff_parts={K2 S3 I0 C1 F1}, diff_note={concepts: Gauss law with non-uniform density, field energy | steps: enclosed charge by integrating rho; E(r) inside and its maximum at r = 20 cm; total charge 0.0225 uC and energy outside depends only on Q; field outside is inverse-square not hyperbola; test each option}]
A non-conducting solid sphere of radius 30cm and relative permittivity 1 has the volume charge density $\rho=\left(\dfrac{5}{2\pi}\,\mu\text{C/m}^3\right)\left(1-\dfrac{r}{30\text{cm}}\right)$ where $r$ (in cm) is the distance of any point in the space from the centre of the sphere. Then choose the correct statement(s).
\begin{qoptions}
\opt{A}{Density of electric field lines is maximum at $r=25$ cm.}
\opt{B}{Density of electric field lines is maximum at $r=20$ cm.}
\opt{C}{Electrical energy stored in the space $r=30$ cm to $r\to\infty$ is same as that of a solid non-conducting sphere carrying charges of 0.0225 $\mu$C and radii 30 cm having uniform volume charge density in the same space $r=30$ cm to $r\to\infty$.}
\opt{D}{Variation of magnitude of electric field due to the sphere with respect to radial distance $r$ is rectangular hyperbola in the region $r=30$ cm to $r\to\infty$.}
\end{qoptions}
\end{bankq}
\begin{bankq}[id=TOAS13-P1-PHY-S3-Q04, source={TOAS - Test 13 (24 Jul 2022)}, paper=1, qno=4, type=numerical, subtopic={Electric field}, answer={10.00}, answer_src=key, difficulty=5, diff_parts={K2 S2 I0 C1 F1}, diff_note={concepts: motion of a charge in a uniform field, potential divider | steps: tan 37 = vy/u gives vy; vy = e V l/(m d u) gives V; capacitor voltage equals E R/(R + r) in steady state; solve R}]
A parallel plate capacitor with plate of length $\ell$ is included in a circuit as shown in fig. Given are, e.m.f. of the current source, its internal resistance $r$ and the distance $d$ between the plates. As electron with a velocity $u$ flies into the capacitor, so that the electron flies out of the capacitor at an angle $37^\circ$ to the plates. Find the value of $R$ in ohm (assuming circuit in steady state and neglect relativistic effect). Take $\ell=91$ cm, $E=3$V, $r=2\Omega$, $d=\dfrac13$ mm, $u=4\times10^7$ m/s, $M_e=9.1\times10^{-31}$ kg and $e=1.6\times10^{-19}$C.
\bfig{116.3}{TOAS13_p1_phy_s3_q04_1.png}
\end{bankq}
\begin{bankq}[id=TOAS13-P2-PHY-S1-Q05, source={TOAS - Test 13 (24 Jul 2022)}, paper=2, qno=5, type=single, subtopic={Electric potential}, answer={C}, answer_src=key, difficulty=2, diff_parts={K1 S1 I1 C0 F0}, diff_note={concepts: potential of a charge system | steps: the net charge is zero so the potential vanishes on the axis through the centre}]
Four charges of 8$\mu$C, $-6\mu$C, 2 $\mu$C and $-4\mu$C are placed at corners of a square placed in xy plane with its center at origin. The electric potential due do these charges is zero every where on the line
\begin{qoptions}
\opt{A}{$x=y$, $z=0$}
\opt{B}{$x=0=z$}
\opt{C}{$x=0=y$}
\opt{D}{$x=z$, $y=0$}
\end{qoptions}
\end{bankq}
\begin{bankq}[id=TOAS13-P2-PHY-S1-Q20, source={TOAS - Test 13 (24 Jul 2022)}, paper=2, qno=20, type=single, subtopic={Electric potential}, answer={C}, answer_src=key, difficulty=2, diff_parts={K1 S1 I1 C0 F0}, diff_note={concepts: energy conservation in a Coulomb field | steps: use the given relation for KE; 1/r = 1/d + 4/d}]
A charged particle of mass `m' and charge $q$ is moving directly towards a fixed charge $Q$. When it is at a distance `d' from fixed charge it has velocity $v$. The distance of closest approach will be. (Here, $\dfrac{\pi\epsilon_0mv^2}{qQ}=\dfrac2d$)
\begin{qoptions}
\opt{A}{$\dfrac d2$}
\opt{B}{$\dfrac d3$}
\opt{C}{$\dfrac d5$}
\opt{D}{$\dfrac{3d}{4}$}
\end{qoptions}
\end{bankq}
\begin{bankq}[id=TOAS13-P2-PHY-S2-Q08, source={TOAS - Test 13 (24 Jul 2022)}, paper=2, qno=8, type=numerical, subtopic={Electric field}, answer={69.00}, answer_src=key, difficulty=2, diff_parts={K1 S1 I0 C0 F1}, diff_note={concepts: electric field from potential | steps: E = -grad V at (2,2,2) is (-40,-40,-40) with magnitude 40 sqrt 3}]
The electric potential existing in space is given as $V(x,y,z)=10\,(xy+yz+zx)$ in SI unit. Find the value of electric field (in N/C) at point $(2,2,2)$ (approximately).
\end{bankq}
\begin{bankq}[id=TOAS15-P1-PHY-S1-Q05, source={TOAS - Test 15 (22 Aug 2022)}, paper=1, qno=5, type=single, subtopic={Electric field}, answer={B}, answer_src=key, difficulty=4, diff_parts={K2 S2 I1 C0 F0}, diff_note={concepts: electrostatic force, stable equilibrium | steps: restoring condition for small displacement; resolve the electric force; minimum q from the stability inequality}]
A small ball of mass $m$, charge $Q$ is to be held in stable equilibrium for small angular displacement as shown in figure. Find minimum charge $q$ that should be fixed below ball for this :-
\bfig{89.9}{TOAS15_p1_phy_s1_q05_1.png}
\begin{qoptions}
\opt{A}{$\dfrac{16\pi\epsilon_0mgR^2}{Q}$}
\opt{B}{$\dfrac{32\pi\epsilon_0mgR^2}{Q}$}
\opt{C}{$\dfrac{8\pi\epsilon_0mgR^2}{3Q}$}
\opt{D}{$\dfrac{64\pi\epsilon_0mgR^2}{Q}$}
\end{qoptions}
\end{bankq}
\begin{bankq}[id=TOAS15-P1-PHY-S2-Q11, source={TOAS - Test 15 (22 Aug 2022)}, paper=1, qno=11, type=multi, subtopic={Electric field}, answer={A,B,C,D}, answer_src=key, difficulty=5, diff_parts={K2 S2 I1 C0 F1}, diff_note={concepts: electric field, potential of a ring segment | steps: charge of an element is proportional to its distance; field varies as 1/r; potential same; net field direction}]
A wire having a uniform linear charge density $\lambda$, is bent in the form of ring of radius $R$. Point A as shown in the figure, is in the plane of ring but not at centre. Two elements of the ring of length $a_1$ and $a_2$ subtend very small same angle at the point A. They are at distance $r_1$ and $r_2$ from the point A respectively.
\bfig{100.1}{TOAS15_p1_phy_s2_q11_1.png}
\begin{qoptions}
\opt{A}{The ratio of charge of elements $a_1$ and $a_2$ is $\dfrac{r_1}{r_2}$.}
\opt{B}{The element $a_1$ produced greater magnitude of electric field at A than $a_2$.}
\opt{C}{The element $a_1$ and $a_2$ produced same potential at A.}
\opt{D}{The direction of net electric field at A is towards element $a_2$.}
\end{qoptions}
\end{bankq}
\begin{bankq}[id=TOAS15-P2-PHY-S1-Q05, source={TOAS - Test 15 (22 Aug 2022)}, paper=2, qno=5, type=multi, subtopic={Electric field}, answer={A,B}, answer_src=solved, working={Field points toward both charges near them (to $-x$ near $q_1$, to $+x$ near $q_2$): both negative. Zero at $2a/3$: $|q_1|/(2a/3)^2=|q_2|/(a/3)^2$ gives $|q_1/q_2|=4$}, difficulty=3, diff_parts={K1 S1 I1 C0 F1}, diff_note={concepts: electric field plot and superposition | steps: signs from the field direction near each charge; null point condition gives the ratio}]
If the electric field $E$ is plotted, with distance from $q_1$ along the line joining of two, then ($E$ is positive along +ve x-axis) it looks as shown in figure. From the plot we can say that
\bfig{128.6}{TOAS15_p2_phy_s1_q05_1.png}
\begin{qoptions}
\opt{A}{$q_1$ and $q_2$ both are negative}
\opt{B}{$\left|\dfrac{q_1}{q_2}\right|=4$}
\opt{C}{$q_1$ is positive and $q_2$ is negative}
\opt{D}{$\left|\dfrac{q_1}{q_2}\right|=\dfrac14$}
\end{qoptions}
\end{bankq}
\begin{bankq}[id=TOAS15-P2-PHY-S3-Q01, source={TOAS - Test 15 (22 Aug 2022)}, paper=2, qno=1, type=numerical, subtopic={Conductors}, answer={1}, answer_src=key, difficulty=2, diff_parts={K1 S1 I0 C0 F1}, diff_note={concepts: induced charges on earthed plates | steps: induced charges are inversely proportional to the gaps; split the charge of B}]
Three identical, parallel conducting plates A, B and C are placed as shown. Switches $S_1$ and $S_2$ are open and can connect A and C to earth when closed. A charge $Q=-3~\mu$C is given to B. If $S_1$ and $S_2$ both are closed together. Find the amount of charge in $\mu$C on inner surface of plate A.
\bfig{84.2}{TOAS15_p2_phy_s3_q01_1.png}
\end{bankq}
\begin{bankq}[id=TOAS16-P1-PHY-S1-Q01, source={TOAS - Test 16 (04 Sep 2022)}, paper=1, qno=1, type=single, subtopic={Electric field}, answer={B}, answer_src=key, difficulty=5, diff_parts={K2 S2 I1 C1 F0}, diff_note={concepts: motion of a charge in a time-varying field, integration of kinematics | steps: a = e k t/(m l); velocity and position by integration; time to cross the gap; speed at that time}]
At $t=0$, an electron $(e,m)$ leaves one plate of a parallel plate capacitor with negligible initial velocity. An accelerating voltage, varying as $V=kt$, is applied between the plates. The separation between the plates is $\ell$, velocity with which the electron reaches the opposite plate is :
\begin{qoptions}
\opt{A}{$\dfrac{ek}{m\ell}\left[\dfrac{3m\ell^2}{ek}\right]^{2/3}$}
\opt{B}{$\dfrac12\dfrac{ek}{m\ell}\left[\dfrac{6m\ell^2}{ek}\right]^{2/3}$}
\opt{C}{$\dfrac{2ek}{m\ell}\left[\dfrac{6m\ell^2}{ek}\right]^{2/3}$}
\opt{D}{$\dfrac{ek}{m\ell}\left[\dfrac{6m\ell^2}{ek}\right]^{1/3}$}
\end{qoptions}
\end{bankq}
\begin{bankq}[id=TOAS16-P1-PHY-S1-Q02, source={TOAS - Test 16 (04 Sep 2022)}, paper=1, qno=2, type=single, subtopic={Electric potential}, answer={B}, answer_src=key, difficulty=1, diff_parts={K1 S1 I0 C0 F0}, diff_note={concepts: electric field as the gradient of potential | steps: E = -grad V at the point; project on the unit vector of A}]
The potential of a certain electric field in the 3-D space is given as $V=\alpha(xy-z^2)$, $\alpha$ is a constant. The projection of $\vec E$ on to the direction of vector $\vec A=\hat i+\hat j+3\hat k$ at a point $(2,1-3)$ is:
\begin{qoptions}
\opt{A}{Zero}
\opt{B}{$\dfrac{-21\alpha}{\sqrt{11}}$}
\opt{C}{$-\sqrt{11}\,\alpha$}
\opt{D}{$\dfrac{1}{\sqrt{11}}$}
\end{qoptions}
\end{bankq}
\begin{bankpassage}[id=TOAS16-P1-PHY-S3-G1-2, source={TOAS - Test 16 (04 Sep 2022)}, questions={1-2}]
A safety engineer needs to evaluate the practice of storing inflammable conducting liquids in non-conducting containers. The company supplying the liquid uses a cylindrical plastic container of radius $r=20$ cm and fills it to a height of $h=20$ cm, which is not equal to its full height. During handling, the exterior surface acquires a charge density of $\sigma=2\times10^{-6}$ C/m$^2$, approximately uniform. This induces a charge separation within the liquid, because it is conducting (see figure).
\bfig{109.4}{TOAS16_p1_phy_s3_para1_2_1.png}
\end{bankpassage}
\begin{bankq}[id=TOAS16-P1-PHY-S3-Q01, source={TOAS - Test 16 (04 Sep 2022)}, paper=1, qno=1, type=numerical, subtopic={Conductors}, answer={0.72}, answer_src=key, difficulty=4, diff_parts={K2 S1 I1 C0 F1}, diff_note={concepts: induced charge on a conducting liquid, surface charge | steps: area of the wetted wall and base; total induced charge equals the opposite surface charge on them}, group=TOAS16-P1-PHY-S3-G1-2]
How much charge is induced within the bulk of the liquid (in $\mu$C) Take : $\pi\approx3$.
\end{bankq}
\begin{bankq}[id=TOAS16-P1-PHY-S3-Q02, source={TOAS - Test 16 (04 Sep 2022)}, paper=1, qno=2, type=numerical, subtopic={Conductors}, answer={7.20}, answer_src=key, difficulty=2, diff_parts={K1 S1 I0 C0 F1}, diff_note={concepts: energy stored in a capacitor | steps: U = Q\^{}2/2C}, group=TOAS16-P1-PHY-S3-G1-2]
Assume that the capacitance of the central part relative to ground is nearly 36pF. Calculate the potential energy stored in the capacitor (in mJ).
\end{bankq}
\begin{bankq}[id=TOAS17-P2-PHY-S2-Q08, source={TOAS - Test 17 (18 Sep 2022)}, paper=2, qno=8, type=numerical, subtopic={Electric potential}, answer={30}, answer_src=key, difficulty=2, diff_parts={K1 S1 I0 C0 F1}, diff_note={concepts: potential from the field graph | steps: potential difference as the area under the graph; subtract}]
A graph of the $x$ component of the electric field as a function of $x$ in a region of space is shown. The $Y$ and $Z$ components of the electric field are zero in this region. If the electric potential is 10V at the origin, then potential (in V) at $x=2.0$ m is :
\bfig{154.2}{TOAS17_p2_phy_s2_q08_1.png}
\end{bankq}
\begin{bankq}[id=TOAS17-P2-PHY-S2-Q09, source={TOAS - Test 17 (18 Sep 2022)}, paper=2, qno=9, type=numerical, subtopic={Gauss's law}, answer={2}, answer_src=key, difficulty=5, diff_parts={K1 S2 I1 C1 F1}, diff_note={concepts: energy density in the field outside a sphere | steps: energy in each region as an integral; equal energies; radius ratio}]
Consider a uniformly charged solid non conducting sphere of radius $R$ and charge $Q$. Consider three spherical regions centered at the centre of the uniformly charged solid non conducting sphere, as given below :\par
Region (1) $R_1\ge r\ge0$\par
Region (2) $R_2\ge r\ge R_1$\par
Region (3) $\infty\ge r\ge R_2$\par
(Both $R_1$ \& $R_2>R$)\par
If electrostatic potential energy stored in each of the three regions is equal then $\dfrac{R_2}{R_1}$ is :
\end{bankq}
\begin{bankq}[id=TOAS18-P1-PHY-S3-Q07, source={TOAS - Test 18 (30 Oct 2022)}, paper=1, qno=7, type=numerical, subtopic={Electric potential}, answer={1.24}, answer_src=solved, working={Zero potential: $|r|^2=4|r+R\hat x|^2$ gives a sphere of radius $2R/3$; volume $\tfrac43\pi(2R/3)^3=\tfrac{32\pi}{81}R^3$, $K\approx1.24$}, difficulty=5, diff_parts={K1 S2 I1 C1 F1}, diff_note={concepts: equipotential surface of two charges | steps: zero potential condition; sphere of Apollonius; volume}]
A particle with charge $q$ is located at $(-R,0,0)$ and a particle with charge $-2q$ is located at origin then equipotential surface that has zero potential occupy $KR^3$ volume in space. Find $K$.
\end{bankq}
\begin{bankq}[id=TOAS18-P2-PHY-S2-Q05, source={TOAS - Test 18 (30 Oct 2022)}, paper=2, qno=5, type=multi, subtopic={Gauss's law}, answer={B,D}, answer_src=key, difficulty=2, diff_parts={K1 S1 I0 C0 F1}, diff_note={concepts: Gauss law | steps: statements about flux with one type of charge}]
If there were only one type of charge in the universe then :
\begin{qoptions}
\opt{A}{$\displaystyle\oint_S\vec E\cdot\overrightarrow{dS}\neq0$ on any surface}
\opt{B}{$\displaystyle\oint_S\vec E\cdot\overrightarrow{dS}=0$ if charge is outside the surface}
\opt{C}{$\displaystyle\oint_S\vec E\cdot d\vec S$ could not be defined}
\opt{D}{$\displaystyle\oint_S\vec E\cdot d\vec S=\dfrac q{\varepsilon_0}$ if charges of magnitude $q$ were inside the surface.}
\end{qoptions}
\end{bankq}
\begin{bankq}[id=TOAS18-P2-PHY-S2-Q06, source={TOAS - Test 18 (30 Oct 2022)}, paper=2, qno=6, type=multi, subtopic={Electric field}, answer={A,B,D}, answer_src=key, difficulty=4, diff_parts={K1 S2 I1 C0 F1}, diff_note={concepts: stability of a charge in a ring field | steps: force in the plane and along the axis; stability; small oscillation period}]
A positive charge $Q$ is uniformly distributed along a non conducting circular ring of radius $R$. A small test charge $q$ of mass $m$ kept on centre of ring. Then :
\bfig{93.7}{TOAS18_p2_phy_s2_q06_1.png}
\begin{qoptions}
\opt{A}{If $q>0$ \& is displaced away from center in plane of ring it will be pushed towards center.}
\opt{B}{If $q<0$ \& is displaced away from center in plane of ring then it will never return to center.}
\opt{C}{If $q<0$ \& is displaced along axis then it will perform SHM.}
\opt{D}{If $q<0$ \& is slightly displaced along axis then time period of oscillation is $2\pi\sqrt{\dfrac{mR^3}{kQq}}$.}
\end{qoptions}
\end{bankq}

\section*{TAAS papers}
\begin{bankq}[id=TAAS01-P1-PHY-S1-Q04, source={TAAS - Test 01 (18 Jun 2023)}, paper=1, qno=4, type=single, subtopic={Electric field}, answer={A}, answer_src=key, difficulty=2, diff_parts={K1 S1 I1 C0 F0}, diff_note={concepts: field of charged wires on edges of a cube | steps: field of each wire at the centre by symmetry; vector sum}]
Three identical thin uniform charged wires are fixed along three sides of a cube as shown in figure. Length of each wire is $l$ and linear charge density on each of them is $\lambda$. Electric field at the center of the cube is :
\bfig{70.1}{TAAS01_p1_phy_s1_q04_1.png}
\begin{qoptions}
\opt{A}{0}
\opt{B}{$\dfrac{3\sqrt2\,\lambda}{4\pi\varepsilon_0l}$}
\opt{C}{$\dfrac{12\sqrt2\,\lambda}{4\pi\varepsilon_0l}$}
\opt{D}{$\dfrac{3\sqrt3\,\lambda}{16\pi\varepsilon_0l}$}
\end{qoptions}
\end{bankq}
\begin{bankq}[id=TAAS01-P1-PHY-S3-Q04, source={TAAS - Test 01 (18 Jun 2023)}, paper=1, qno=4, type=numerical, subtopic={Gauss's law}, answer={0.35 to 0.40}, answer_src=key, difficulty=7, diff_parts={K1 S3 I2 C1 F1}, diff_note={concepts: field of overlapping uniformly charged spheres | steps: superposition of three uniform spheres; centres; field inside each; vector sum; value of N}]
In the figure shown, sphere $S_1$, $S_1$ and $S_3$ have radii $R$, $\dfrac R2$ \& $\dfrac R4$ $C_1$, $C_2$ \& $C_3$ are centers of the three spheres lying in a plane. Angle $\angle C_1C_2C_3$ is right angle. Sphere $S_3$ has a uniformly spread volume charge density $4\rho$. The remaining part of $S_2$ has uniform charge density of $2\rho$ and the left over part of $S_1$ has a uniform charge density of $\rho$. The electric field at a point $A$ at a distance $\dfrac R8$ from $C_3$ on the line $C_2C_3$ is $\dfrac{N\rho R}{\varepsilon_0}$. Find $N$
\bfig{100.6}{TAAS01_p1_phy_s3_q04_1.png}
\end{bankq}
\begin{bankq}[id=TAAS02-P1-PHY-S1-Q01, source={TAAS - Test 02 (30 Jul 2023)}, paper=1, qno=1, type=single, subtopic={Electric potential}, answer={A}, answer_src=key, difficulty=4, diff_parts={K1 S2 I1 C1 F0}, diff_note={concepts: work in the field of a line charge | steps: potential of a line charge by logarithm; distances of C and A; work q times the potential difference}]
For an infinite line of charge having charge density $\lambda$ lying along $x$-axis, the work required in moving charge $q$ from $C$ to $A$ along arc $CA$ is :-
\bfig{126.9}{TAAS02_p1_phy_s1_q01_1.png}
\begin{qoptions}
\opt{A}{$\dfrac{q\lambda}{\pi\varepsilon_0}\log_e\sqrt2$}
\opt{B}{$\dfrac{q\lambda}{4\pi\varepsilon_0}\log_e\sqrt2$}
\opt{C}{$\dfrac{q\lambda}{4\pi\varepsilon_0}\log_e2$}
\opt{D}{$\dfrac{q\lambda}{2\pi\varepsilon_0}\log_e\dfrac12$}
\end{qoptions}
\end{bankq}
\begin{bankq}[id=TAAS02-P1-PHY-S4-Q06, source={TAAS - Test 02 (30 Jul 2023)}, paper=1, qno=6, type=numerical, subtopic={Electric field}, answer={6}, answer_src=key, difficulty=7, diff_parts={K1 S3 I2 C1 F1}, diff_note={concepts: field of a charged sphere with a cavity | steps: uniform field inside the cavity; force on the electron; acceleration components; time to reach the surface; P}]
A cavity of radius $r$ is present inside a fixed solid dielectric sphere of radius $R$, having a volume charge density of $\rho$. The distance between the centres of the sphere and the cavity is $a$. An electron $e$ is kept inside the cavity at an angle $\theta=45^\circ$ as shown. The electron (of mass $m$ and charge $-e$) will take $\left(\dfrac{P\sqrt2mr\varepsilon_0}{ea\rho}\right)^{1/2}$ time to touch the sphere again. Neglect gravity. Find the value of $P$.
\bfig{79.4}{TAAS02_p1_phy_s4_q06_1.png}
\end{bankq}
\begin{bankq}[id=TAAS02-P2-PHY-S4-Q05, source={TAAS - Test 02 (30 Jul 2023)}, paper=2, qno=5, type=numerical, subtopic={Conductors}, answer={2}, answer_src=key, difficulty=6, diff_parts={K1 S2 I2 C1 F1}, diff_note={concepts: connected conducting spheres inside a grounded shell | steps: capacitances; potential after placing inside the shell; charges; charge flowed}]
Two metal balls of radii $R_1$ and $R_2$ such that $R_2=2R_1$ is connected to a long thin conducting wire and has total charge $Q=21\mu$C. Then the ball of radius $R_1$ is placed inside a grounded metal sphere of radius $R=3R_1$ as shown. What charge (in $\mu$C) will flow through the wire in this process ?
\bfig{108.3}{TAAS02_p2_phy_s4_q05_1.png}
\end{bankq}
\begin{bankq}[id=TAAS02-P2-PHY-S4-Q06, source={TAAS - Test 02 (30 Jul 2023)}, paper=2, qno=6, type=numerical, subtopic={Gauss's law}, answer={6}, answer_src=key, difficulty=6, diff_parts={K1 S2 I2 C1 F1}, diff_note={concepts: flux of a ring field through a sphere | steps: field of a ring along its axis; the sphere position; Gauss law; maximum flux}]
A ring of radius $R$ having a linear charge density $\lambda$ moves towards a solid imaginary sphere of radius $\dfrac R2$, so that the centre of ring passes through the centre of sphere. The axis of the ring is perpendicular to the line joining the centres of the ring and the sphere. If the maximum flux through the sphere in this process is $\dfrac{2\lambda\pi R}{n\in_0}$ then find the value of $n$?
\bfig{86.2}{TAAS02_p2_phy_s4_q06_1.png}
\end{bankq}

\section*{TMAS papers}
\begin{bankq}[id=TMAS01-P1-PHY-S1-Q01, source={TMAS - Test 01 (09 Jul 2023)}, paper=1, qno=1, type=single, subtopic={Gauss's law}, answer={C}, answer_src=key, difficulty=3, diff_parts={K1 S1 I1 C1 F0}, diff_note={concepts: flux through a closed surface | steps: divergence of the field; volume integral}]
Electric field at point $P$ is given by $\vec E=E_0r\,\hat r$. The total flux through the given cylinder of radius $R$ and height $h$ is :-
\bfig{55.6}{TMAS01_p1_phy_s1_q01_1.png}
\begin{qoptions}
\opt{A}{$E_0\pi R^2h$}
\opt{B}{$2E_0\pi R^2h$}
\opt{C}{$3E_0\pi R^2h$}
\opt{D}{$4E_0\pi R^2h$}
\end{qoptions}
\end{bankq}
\begin{bankq}[id=TMAS01-P1-PHY-S1-Q02, source={TMAS - Test 01 (09 Jul 2023)}, paper=1, qno=2, type=single, subtopic={Electric potential}, answer={A}, answer_src=key, difficulty=2, diff_parts={K1 S1 I0 C1 F0}, diff_note={concepts: equipotential surface | steps: work on an equipotential; field lines perpendicular}]
\textbf{Statement-1 :} Work done by electric field on moving a positive charge on an equipotential surface is always zero.\\
and\\
\textbf{Statement-2 :} Electric lines of forces are perpendicular to equipotential surface.
\begin{qoptions}
\opt{A}{Statement--1 is True, Statement--2 is True ; Statement--2 is a correct explanation for Statement--1.}
\opt{B}{Statement--1 is True, Statement--2 is True ; Statement--2 is not a correct explanation for Statement--1.}
\opt{C}{Statement--1 is True, Statement--2 is False.}
\opt{D}{Statement--1 is False, Statement--2 is True.}
\end{qoptions}
\end{bankq}
\begin{bankq}[id=TMAS01-P1-PHY-S1-Q03, source={TMAS - Test 01 (09 Jul 2023)}, paper=1, qno=3, type=single, subtopic={Electric field}, answer={D}, answer_src=key, difficulty=3, diff_parts={K1 S1 I1 C1 F0}, diff_note={concepts: vector sum of fields | steps: forces from each corner; direction of the sum}]
Four charges are arranged at the corners of a square $ABCD$, as shown. The net force on positive charge kept at the centre of the square is
\bfig{63.4}{TMAS01_p1_phy_s1_q03_1.png}
\begin{qoptions}
\opt{A}{zero}
\opt{B}{along diagonal $AC$}
\opt{C}{along diagonal $BD$}
\opt{D}{perpendicular to the side $AB$}
\end{qoptions}
\end{bankq}
\begin{bankq}[id=TMAS01-P1-PHY-S1-Q04, source={TMAS - Test 01 (09 Jul 2023)}, paper=1, qno=4, type=single, subtopic={Conductors}, answer={C}, answer_src=key, difficulty=3, diff_parts={K1 S1 I1 C1 F0}, diff_note={concepts: field near a conductor | steps: field of surface charge; external charges contribution}]
The density of charge at point $P$ on the conductor is $\sigma$. The resultant electric field near $P$ (just outside) will be :-
\bfig{166.3}{TMAS01_p1_phy_s1_q04_1.png}
\begin{qoptions}
\opt{A}{$\dfrac{kQ_1}{r_1^2}\left(-\hat j\right)+\dfrac{kQ_2}{r_2^2}\left(-\hat i\right)$}
\opt{B}{$\dfrac{kQ_1}{r_1^2}\left(-\hat j\right)+\dfrac{kQ_2}{r_2^2}\left(-\hat i\right)+\dfrac\sigma{\in_0}\left(\hat i\right)$}
\opt{C}{$\dfrac\sigma{\in_0}\hat i$}
\opt{D}{$\dfrac{kQ_1}{r_1^2}\left(\hat j\right)+\dfrac{kQ_2}{r_2^2}\left(\hat i\right)+\dfrac\sigma{\in_0}\left(-\hat i\right)$}
\end{qoptions}
\end{bankq}
\begin{bankq}[id=TMAS01-P1-PHY-S1-Q05, source={TMAS - Test 01 (09 Jul 2023)}, paper=1, qno=5, type=single, subtopic={Gauss's law}, answer={A}, answer_src=key, difficulty=2, diff_parts={K1 S1 I0 C1 F0}, diff_note={concepts: Gauss law | steps: enclosed charges; sign}]
Find out electric flux passing from a closed Gaussian surface as shown in figure :
\bfig{137.6}{TMAS01_p1_phy_s1_q05_1.png}
\begin{qoptions}
\opt{A}{$\dfrac{q_1+q_5-q_6}{\varepsilon_0}$}
\opt{B}{$\dfrac{q_1+q_5+q_6}{\varepsilon_0}$}
\opt{C}{$\dfrac{q_1+q_2+q_3+q_4+q_5+q_6}{\varepsilon_0}$}
\opt{D}{$\dfrac{q_1+q_2-q_3+q_4+q_5-q_6}{\varepsilon_0}$}
\end{qoptions}
\end{bankq}
\begin{bankq}[id=TMAS01-P1-PHY-S1-Q06, source={TMAS - Test 01 (09 Jul 2023)}, paper=1, qno=6, type=single, subtopic={Electric potential}, answer={C}, answer_src=key, difficulty=4, diff_parts={K1 S2 I1 C1 F0}, diff_note={concepts: potential graph | steps: sign from the graph; zero field point; magnitudes}]
Two point charges $Q_1$ and $Q_2$ lie along a line at a distance $a$ from each other. Figure shows the potential variation along line of charges then which of the following statement is incorrect?
\bfig{165.3}{TMAS01_p1_phy_s1_q06_1.png}
\begin{qoptions}
\opt{A}{Electric field is zero at point 3}
\opt{B}{$Q_1$ is positive and $Q_2$ is negative}
\opt{C}{$Q_1$ is negative and $Q_2$ is positive}
\opt{D}{$|Q_1|>|Q_2|$}
\end{qoptions}
\end{bankq}
\begin{bankq}[id=TMAS01-P1-PHY-S1-Q18, source={TMAS - Test 01 (09 Jul 2023)}, paper=1, qno=18, type=single, subtopic={Electric potential}, answer={A}, answer_src=key, difficulty=3, diff_parts={K1 S1 I1 C1 F0}, diff_note={concepts: potential from charge distribution | steps: charges on the ring; sum of potentials}]
Charge distribution on a ring in $x$-$y$ plane is shown in the figure. Then electric potential at $O$ is :
\bfig{103.1}{TMAS01_p1_phy_s1_q18_1.png}
\begin{qoptions}
\opt{A}{$\dfrac Q{4\pi\in_0r}$}
\opt{B}{Zero}
\opt{C}{$\dfrac Q{2\pi\in_0r}$}
\opt{D}{$\dfrac{3Q}{4\pi\in_0r}$}
\end{qoptions}
\end{bankq}
\begin{bankq}[id=TMAS01-P1-PHY-S1-Q19, source={TMAS - Test 01 (09 Jul 2023)}, paper=1, qno=19, type=single, subtopic={Electric potential}, answer={D}, answer_src=key, difficulty=3, diff_parts={K1 S1 I1 C1 F0}, diff_note={concepts: equipotential and field | steps: field perpendicular to the line; check}]
The equation of an equipotential line in an electric field is $y=2x$, then the electric field strength vector at $(1,2)$ may be :
\begin{qoptions}
\opt{A}{$4\hat i+3\hat j$}
\opt{B}{$4\hat i+8\hat j$}
\opt{C}{$8\hat i+4\hat j$}
\opt{D}{$-8\hat i+4\hat j$}
\end{qoptions}
\end{bankq}
\begin{bankq}[id=TMAS01-P1-PHY-S1-Q20, source={TMAS - Test 01 (09 Jul 2023)}, paper=1, qno=20, type=single, subtopic={Coulomb's law}, answer={D}, answer_src=key, difficulty=5, diff_parts={K2 S2 I1 C1 F0}, diff_note={concepts: charge sharing, Coulomb force | steps: force before; charges after joining; force after}]
Two identical metal balls with charge $+2Q$ and $-Q$ are separated by some distance, and exert a force $F$ on each other. They are joined by a long conducting wire, which is then removed. The force between them will now be
\begin{qoptions}
\opt{A}{$F$}
\opt{B}{$F/2$}
\opt{C}{$F/4$}
\opt{D}{$F/8$}
\end{qoptions}
\end{bankq}
\begin{bankq}[id=TMAS01-P1-PHY-S2-Q01, source={TMAS - Test 01 (09 Jul 2023)}, paper=1, qno=1, type=numerical, subtopic={Gauss's law}, answer={56}, answer_src=key, difficulty=6, diff_parts={K2 S2 I1 C1 F1}, diff_note={concepts: Gauss law, flux through faces | steps: flux through each face; non uniform face; charge enclosed}]
An electric field is given by $\vec E=5\hat i-(7y+8)\hat j-3\hat k$ (S.I. unit) crosses gaussian cube having side length 2m as shown in the figure. Then the net charge enclosed in the cube in $-x\in_0$ (S.I. unit) then find $x$.
\bfig{86.9}{TMAS01_p1_phy_s2_q01_1.png}
\end{bankq}
\begin{bankq}[id=TMAS01-P1-PHY-S2-Q02, source={TMAS - Test 01 (09 Jul 2023)}, paper=1, qno=2, type=numerical, subtopic={Electric potential}, answer={32}, answer_src=key, difficulty=5, diff_parts={K1 S2 I1 C1 F1}, diff_note={concepts: work to assemble charges | steps: potential at the centre; work equals charge times potential; form}]
Four $+q$ charges are fixed on corner of square of side $a$. A fifth charge $+Q$ is brought by external agent slowly from infinity and placed at centre of square. Work done by external agent is $\dfrac{\sqrt\alpha kQq}a$ (where $k=\dfrac1{4\pi\in_0}$). Fill $\alpha$.
\end{bankq}
\begin{bankq}[id=TMAS01-P1-PHY-S2-Q03, source={TMAS - Test 01 (09 Jul 2023)}, paper=1, qno=3, type=numerical, subtopic={Electric field}, answer={2}, answer_src=key, difficulty=8, diff_parts={K2 S2 I2 C2 F1}, diff_note={concepts: field line geometry, flux conservation | steps: flux through a cone; relation of the angles; solve}]
Two charges $+q_1$ and $-q_2$ are placed at $A$ and $B$ respectively. A line of force emanates from $q_1$ at an angle $\alpha$ with the line $AB$. If it terminates at $-q_2$ at an angle $P\sin^{-1}\left[\sqrt{\dfrac{q_1}{q_2}}\sin\dfrac\alpha2\right]$. Find the value of $P$
\end{bankq}
\begin{bankq}[id=TMAS01-P1-PHY-S2-Q04, source={TMAS - Test 01 (09 Jul 2023)}, paper=1, qno=4, type=numerical, subtopic={Electric field}, answer={9}, answer_src=key, difficulty=6, diff_parts={K1 S2 I2 C1 F1}, diff_note={concepts: force on a line segment | steps: field of an infinite wire; integrate along the segment; evaluate}]
An infinitely long wire with linear charge density $\lambda_1$ and a segment of length $\ell$ with charge density $\lambda_2$ lie in a plane at right angle to each other and seperated by a distance $r_0$ as shown in figure. The force with which these two interact is $\dfrac{\lambda_1\lambda_2}{4\pi\varepsilon_0}\,\ell n\,x$. If $\ell=2r_0$, then find the value of $x$.
\bfig{145.7}{TMAS01_p1_phy_s2_q04_1.png}
\end{bankq}
\begin{bankq}[id=TMAS01-P1-PHY-S2-Q05, source={TMAS - Test 01 (09 Jul 2023)}, paper=1, qno=5, type=numerical, subtopic={Gauss's law}, answer={50}, answer_src=key, difficulty=3, diff_parts={K1 S1 I0 C1 F1}, diff_note={concepts: electric flux | steps: dot product with the area}]
Find out the electric flux (S.I. unit) through an area 10 m$^2$ lying in $XY$ plane due to an electric field $\vec E=2\hat i-10\hat j+5\hat k$ (S.I. unit).
\end{bankq}
\begin{bankq}[id=TMAS01-P2-PHY-S1-Q01, source={TMAS - Test 01 (09 Jul 2023)}, paper=2, qno=1, type=multi, subtopic={Gauss's law}, answer={B,C}, answer_src=key, difficulty=7, diff_parts={K2 S2 I2 C1 F1}, diff_note={concepts: Gauss law for a slab, potential graph | steps: field from the slope; charge density from the second derivative; check units}]
The figure shows the variation of electrostatic potential $V$ in volt with the distance of position of point along $x$ axis from origin due to continuous volume charge distribution in the region $x=-1$m to $x=+1$m, the graph is parabolic ($V=15-5x^2$) and rest portion of graph is straight line. Choose the CORRECT option(s)\\
($\varepsilon_0\to$ permittivity of free space $=8.85\times10^{-12}\ \mathrm{N^{-1}m^{-2}C^2}$). The direction of $\vec E$ along positive $x$-axis is considered as positive.
\bfig{184.1}{TMAS01_p2_phy_s1_q01_1.png}
\begin{qoptions}
\opt{A}{This graph of potential may be due to a thick sheet of infinite dimension ($-1\mathrm m\le x\le1$, $-\infty<y<\infty$ and $-\infty<z<\infty$) with constant volume charge density $1.77\times10^{-19}$ C/m$^2$.}
\opt{B}{This graph of potential may be due to a thick sheet of infinite dimension ($-1\mathrm m\le x\le1$, $-\infty<y<\infty$ and $-\infty<z<\infty$) with constant volume charge density $0.885\times10^{-10}$ C/m$^3$.}
\opt{C}{\bfigi{103.9}{TMAS01_p2_phy_s1_q01_C.png}}
\opt{D}{\bfigi{104.4}{TMAS01_p2_phy_s1_q01_D.png}}
\end{qoptions}
\end{bankq}
\begin{bankq}[id=TMAS01-P2-PHY-S1-Q02, source={TMAS - Test 01 (09 Jul 2023)}, paper=2, qno=2, type=multi, subtopic={Dipoles}, answer={A,D}, answer_src=key, difficulty=8, diff_parts={K2 S2 I2 C2 F1}, diff_note={concepts: interaction of ring dipoles, charge distribution | steps: dipole moment of each ring; interaction energy; force; check}]
The given two fixed rings of radius $R$ lie in $XY$ plane as shown in the figure. Linear charge density of the rings $A$ and $B$ varies as per the relation $\lambda=\lambda_0\sin\theta$ and $\lambda=\lambda_0\cos\theta$ respectively where $\theta$ is measured from $+x$ axis. Distance between the centre of the two rings is $r\gg R$.
\bfig{126.9}{TMAS01_p2_phy_s1_q02_1.png}
\begin{qoptions}
\opt{A}{Force of interaction between the two ring is $\dfrac34\times\dfrac{\lambda_0^2\pi}{\in_0}\left(\dfrac Rr\right)^4$.}
\opt{B}{Force of interaction between the two ring is $\dfrac32\times\dfrac{\lambda_0^2\pi}{\in_0}\left(\dfrac Rr\right)^4$.}
\opt{C}{Potential energy of interaction between the two ring is $\dfrac{\pi\lambda_0^2R^4}{4\varepsilon_0r^3}$.}
\opt{D}{Potential energy of interaction between the two ring is zero.}
\end{qoptions}
\end{bankq}
\begin{bankq}[id=TMAS01-P2-PHY-S1-Q03, source={TMAS - Test 01 (09 Jul 2023)}, paper=2, qno=3, type=multi, subtopic={Gauss's law}, answer={A,B,C}, answer_src=key, difficulty=6, diff_parts={K1 S2 I2 C1 F1}, diff_note={concepts: flux through parts of a cube | steps: symmetry; fraction of the flux; compare}]
A point charge $q$ is placed at $P(0,0,a)$. Then select correct alternatives.
\bfig{129.2}{TMAS01_p2_phy_s1_q03_1.png}
\begin{qoptions}
\opt{A}{$\dfrac q{24\varepsilon_0}>\phi_{OAB}>\dfrac q{48\varepsilon_0}$}
\opt{B}{$\phi_{OAC}=\dfrac q{48\varepsilon_0}$}
\opt{C}{$\phi_{ABC}<\phi_{ABO}$}
\opt{D}{$\phi_{OAC}<\phi_{OBC}$}
\end{qoptions}
\end{bankq}
\begin{bankpassage}[id=TMAS01-P2-PHY-S2-G8-10, source={TMAS - Test 01 (09 Jul 2023)}, questions={8-10}]
Three columns are given. Column--I shows conductor in different situation, whereas column--II and column-III represents magnitude of electric field and modulus of potential at point `$P$' respectively.

\noindent\begin{minipage}{\textwidth}
\footnotesize\setlength{\tabcolsep}{4pt}
\begin{tabular}{|l|p{6.6cm}|l|l|l|l|}
\hline
\multicolumn{2}{|c|}{Column-I} & \multicolumn{2}{c|}{Column-II} & \multicolumn{2}{c|}{Column-III}\\
\hline
(I) & Electric field \& potential at `$P$' due to charges on shell is $E$ \& $V$ respectively.\par\bfigi{80.4}{TMAS01_p2_phy_s2_para8_10_1.png} & (i) & $E=\dfrac{5Kq}{8R^2}$ & (P) & $V=\dfrac{15Kq}{6R}$\\
\hline
(II) & Solid conductor of inner radius $R$ and outer radius $2R$. Net Electric field \& potential at `$P$' is $E$ \& $V$ respectively.\par\bfigi{56.9}{TMAS01_p2_phy_s2_para8_10_2.png} & (ii) & $E=\dfrac{4Kq}{9R^2}$ & (Q) & $V=\dfrac{4Kq}{6R}$\\
\hline
(III) & Electric field \& potential at `$P$' due to charges on shell is $E$ \& $V$ respectively.\par\bfigi{92.2}{TMAS01_p2_phy_s2_para8_10_3.png} & (iii) & $E=\dfrac{2Kq}{9R^2}$ & (R) & $V=\dfrac{17Kq}{6R}$\\
\hline
(IV) & $P$ point lies on the axis of co-axial rings having radius $R$ and $\sqrt3R$ \& charges $\sqrt2q$ \& $q$ respectively.\par\bfigi{134.7}{TMAS01_p2_phy_s2_para8_10_4.png} & (iv) & $E=\dfrac{4Kq}{R^2}$ & (S) & $V=\dfrac{9Kq}{6R}$\\
\hline
\end{tabular}
\end{minipage}
\end{bankpassage}
\begin{bankq}[id=TMAS01-P2-PHY-S2-Q08, source={TMAS - Test 01 (09 Jul 2023)}, paper=2, qno=8, type=single, subtopic={Conductors}, answer={D}, answer_src=key, difficulty=8, diff_parts={K3 S2 I2 C2 F0}, diff_note={concepts: conductors, shells, potentials | steps: fields for each case; potentials for each case; match; check}, group=TMAS01-P2-PHY-S2-G8-10]
Choose the correct option:
\begin{qoptions}
\opt{A}{(IV)-(i)-(Q)}
\opt{B}{(III)-(iv)-(P)}
\opt{C}{(III)-(iii)-(Q)}
\opt{D}{None}
\end{qoptions}
\end{bankq}
\begin{bankq}[id=TMAS01-P2-PHY-S2-Q09, source={TMAS - Test 01 (09 Jul 2023)}, paper=2, qno=9, type=single, subtopic={Conductors}, answer={B}, answer_src=key, difficulty=8, diff_parts={K3 S2 I2 C2 F0}, diff_note={concepts: conductors, shells, potentials | steps: fields for each case; potentials for each case; match; check}, group=TMAS01-P2-PHY-S2-G8-10]
Choose the correct option:
\begin{qoptions}
\opt{A}{(II)-(i)-(Q)}
\opt{B}{(I)-(ii)-(R)}
\opt{C}{(I)-(iii)-(Q)}
\opt{D}{None}
\end{qoptions}
\end{bankq}
\begin{bankq}[id=TMAS01-P2-PHY-S2-Q10, source={TMAS - Test 01 (09 Jul 2023)}, paper=2, qno=10, type=single, subtopic={Conductors}, answer={A}, answer_src=key, difficulty=8, diff_parts={K3 S2 I2 C2 F0}, diff_note={concepts: conductors, shells, potentials | steps: fields for each case; potentials for each case; match; check}, group=TMAS01-P2-PHY-S2-G8-10]
Choose the correct option:
\begin{qoptions}
\opt{A}{(II)-(iv)-(P)}
\opt{B}{(II)-(iii)-(P)}
\opt{C}{(II)-(iii)-(Q)}
\opt{D}{None}
\end{qoptions}
\end{bankq}
\begin{bankq}[id=TMAS01-P2-PHY-S3-Q05, source={TMAS - Test 01 (09 Jul 2023)}, paper=2, qno=5, type=numerical, subtopic={Electric potential}, answer={5}, answer_src=key, difficulty=7, diff_parts={K2 S2 I2 C1 F1}, diff_note={concepts: superposition of spheres, potential | steps: potential of a full sphere; potential of the cavity sphere; subtract; n}]
A solid sphere of radius `$R$' is uniformly charged with charge density $\rho$ in its volume. A spherical cavity of radius $\dfrac R2$ is made in the sphere as shown in the figure. The electric potential at the centre of the sphere is $\dfrac{n\rho R^2}{12\in_0}$. Find the value of $n$.
\bfig{69.3}{TMAS01_p2_phy_s3_q05_1.png}
\end{bankq}

\section*{TNAS papers}
\begin{bankq}[id=TNAS02-P1-PHY-S2-Q05, source={TNAS - Test 02 (07 May 2023)}, paper=1, qno=5, type=multi, subtopic={Gauss's law}, answer={A,D}, answer_src=key, difficulty=6, diff_parts={K1 S2 I2 C1 F1}, diff_note={concepts: flux through a cube face from a diagonal line charge | steps: symmetry among faces; flux through the faces; value; compare}]
An infinite line charge carrying a uniform line charge of density $\lambda$ is placed along the body diagonal AG of a cube of side $a$ as shown. Then
\bfig{77.4}{TNAS02_p1_phy_s2_q05_1.png}
\begin{qoptions}
\opt{A}{The flux of the electric field of this charge through one of the faces of the cube is $\dfrac{\lambda a}{2\sqrt3\,\varepsilon_0}$.}
\opt{B}{The flux of the electric field of this charge through one of the faces of the cube is $\dfrac{\sqrt3\lambda a}{2\varepsilon_0}$.}
\opt{C}{The flux of the electric field of this charge through one of the faces of the cube is $\dfrac{\lambda a}{4\sqrt3\,\varepsilon_0}$.}
\opt{D}{The flux of the electric field through the face ABFE is same as that through BCGF.}
\end{qoptions}
\end{bankq}
\begin{bankq}[id=TNAS02-P1-PHY-S2-Q07, source={TNAS - Test 02 (07 May 2023)}, paper=1, qno=7, type=multi, subtopic={Electric potential}, answer={B,D}, answer_src=key, difficulty=6, diff_parts={K1 S2 I2 C1 F1}, diff_note={concepts: equipotentials of opposite line charges | steps: potential of each line; condition; circle in the plane; radius}]
Two infinite long thin wire have linear charge density $-\lambda$ and $+\lambda$. They are parallel to $z$-axis passing through $x$-axis at point $x=-a$ and $x=+a$ respectively. The equipotential surface having potential $\dfrac{\lambda\ell n(2)}{4\pi\varepsilon_0}$ is
\begin{qoptions}
\opt{A}{a sphere}
\opt{B}{A cylinder}
\opt{C}{having radius of $\sqrt2a$}
\opt{D}{having radius $2\sqrt2a$}
\end{qoptions}
\end{bankq}
\begin{bankq}[id=TNAS02-P1-PHY-S2-Q08, source={TNAS - Test 02 (07 May 2023)}, paper=1, qno=8, type=multi, subtopic={Conductors}, answer={A,C,D}, answer_src=key, difficulty=6, diff_parts={K1 S2 I2 C1 F1}, diff_note={concepts: conducting sphere with external charge | steps: induced charge distribution; potential of the sphere; centre potential; field in cavity}]
A positive point charge having charge $Q$ is placed at a distance $a/2$ from centre of an uncharged spherical conductor of inner radius $a$ and outer radius $2a$. Another positive point charge having charge $Q$ is placed at distance $4a$ from centre of a conducting sphere. Then which of the following is correct ?
\bfig{111.1}{TNAS02_p1_phy_s2_q08_1.png}
\begin{qoptions}
\opt{A}{The surface charge density on inner and outer surface of sphere would be non uniform.}
\opt{B}{The electrostatic potential on outer surface of conductor is $\dfrac{3KQ}{2a}$.}
\opt{C}{Electric potential at centre of sphere is $\dfrac{7KQ}{4a}$.}
\opt{D}{Electric field inside cavity is non-uniform.}
\end{qoptions}
\end{bankq}
\begin{bankq}[id=TNAS02-P2-PHY-S1-Q03, source={TNAS - Test 02 (07 May 2023)}, paper=2, qno=3, type=multi, subtopic={Conductors}, answer={A,B,D}, answer_src=key, difficulty=6, diff_parts={K1 S2 I2 C1 F1}, diff_note={concepts: conducting shell in an external field | steps: field inside; force components; check each}]
A spherical conducting shell having total charge $10$ C is fixed inside a uniform electric field of $40$ N/C. Net force on $4$C charge placed inside shell at distance $R/2$ from centre is $-60\,\mathrm N\,\hat j$. Choose incorrect option(s).
\bfig{71.5}{TNAS02_p2_phy_s1_q03_1.png}
\begin{qoptions}
\opt{A}{Force on $4$C due to external electric field is zero.}
\opt{B}{Force on sphere due to external electric field is $560$ N.}
\opt{C}{Force on $4$C charge particle due to charge on outer surface of spherical shell is $-160\,\mathrm N\,\hat j$.}
\opt{D}{Force on $4$C charge due to spherical shell is $220\,\mathrm N\,\hat j$.}
\end{qoptions}
\end{bankq}
\begin{bankpassage}[id=TNAS02-P2-PHY-S3-G1-2, source={TNAS - Test 02 (07 May 2023)}, questions={1-2}]
A charge particle $q$ is held between two conducting shells as shown. Both of conducting shell are neutral. (Take $q=2\,\mu$C, $R=20$ m)
\bfig{87.0}{TNAS02_p2_phy_s3_para1_2_1.png}
\end{bankpassage}
\begin{bankq}[id=TNAS02-P2-PHY-S3-Q01, source={TNAS - Test 02 (07 May 2023)}, paper=2, qno=1, type=single, subtopic={Conductors}, answer={525}, answer_src=solved, working={$V_{inner}=kq/3R$ and $V_{outer}=kq/4R$; sum $=\dfrac{7kq}{12R}=525$ V}, difficulty=5, diff_parts={K1 S2 I2 C1 F0}, diff_note={concepts: potentials of two shells with a charge between | steps: induced charges; potential of each; sum}, group=TNAS02-P2-PHY-S3-G1-2]
Sum of potential (in V) of both spheres will be :-
\end{bankq}
\begin{bankq}[id=TNAS02-P2-PHY-S3-Q02, source={TNAS - Test 02 (07 May 2023)}, paper=2, qno=2, type=single, subtopic={Conductors}, answer={0.225}, answer_src=solved, working={Heat equals the electrostatic energy released when $q$ leaves the outer surface at potential $kq/4R$: $\dfrac{kq^2}{8R}=2.25\times10^{-4}$ J $=0.225$ mJ}, difficulty=5, diff_parts={K1 S2 I2 C1 F0}, diff_note={concepts: grounding through a resistor | steps: charge flowing; energy change; heat}, group=TNAS02-P2-PHY-S3-G1-2]
If outer surface is grounded by connecting with a resistance $R$, then heat loss (in mJ) in resistance $R$ is
\end{bankq}
\begin{bankq}[id=TNAS03-P1-PHY-S1-Q01, source={TNAS - Test 03 (28 May 2023)}, paper=1, qno=1, type=single, subtopic={Coulomb's law}, answer={D}, answer_src=key, difficulty=4, diff_parts={K1 S2 I1 C1 F0}, diff_note={concepts: equilibrium of charges in a charged sphere | steps: field inside a uniform sphere; mutual force; balance condition}]
In a uniformly charged sphere of charge $2q$ and radius $R$, two identical particles, each of charge $-q$ are embedded in it. If the system is in static equilibrium, the identical particles should be :-
\begin{qoptions}
\opt{A}{at the ends of a diameter of the sphere.}
\opt{B}{one at the centre, other on the surface of the sphere.}
\opt{C}{at two interior points, separated by a distance $2R$.}
\opt{D}{symmetrically placed about the centre with the separation $R$.}
\end{qoptions}
\end{bankq}
\begin{bankq}[id=TNAS03-P1-PHY-S1-Q10, source={TNAS - Test 03 (28 May 2023)}, paper=1, qno=10, type=single, subtopic={Gauss's law}, answer={C}, answer_src=key, difficulty=5, diff_parts={K1 S2 I2 C1 F0}, diff_note={concepts: neutral system of two charged bodies | steps: charge of the inner sphere by integration; charge of the shell; equate}]
A solid sphere of radius $R_1$ and volume charge density $\rho=\dfrac{\rho_0}{r}$ is enclosed in a hollow sphere of radius $R_2$ with negative surface charge density $\sigma$, such that the total charge in the system is zero $\rho_0$ is a positive constant and $r$ is the distance from the centre of the sphere. The ratio $\dfrac{R_2}{R_1}$ is :-
\begin{qoptions}
\opt{A}{$\dfrac{\sigma}{\rho_0}$}
\opt{B}{$\sqrt{\dfrac{2\sigma}{\rho_0}}$}
\opt{C}{$\sqrt{\dfrac{\rho_0}{2\sigma}}$}
\opt{D}{$\dfrac{\rho_0}{\sigma}$}
\end{qoptions}
\end{bankq}
\begin{bankq}[id=TNAS03-P1-PHY-S1-Q18, source={TNAS - Test 03 (28 May 2023)}, paper=1, qno=18, type=single, subtopic={Dipoles}, answer={A}, answer_src=key, difficulty=4, diff_parts={K1 S2 I1 C1 F0}, diff_note={concepts: torque on a dipole | steps: torque in two fields; compare; angle}]
An electric dipole has a fixed dipole moment $\vec p$, which makes angle $\theta$ with respect to $x$-axis. When subjected to an electric field $\vec E_1=E\hat i$, it experiences a torque $\vec T_1=\tau\hat k$. When subjected to another electric field $\vec E_2=\sqrt3E_1\hat j$ it experiences torque $\vec T_2=-\vec T_1$. The angle $\theta$ is :
\begin{qoptions}
\opt{A}{$60^\circ$}
\opt{B}{$90^\circ$}
\opt{C}{$30^\circ$}
\opt{D}{$45^\circ$}
\end{qoptions}
\end{bankq}
\begin{bankq}[id=TNAS03-P2-PHY-S1-Q02, source={TNAS - Test 03 (28 May 2023)}, paper=2, qno=2, type=multi, subtopic={Electric potential}, answer={A,C}, answer_src=key, difficulty=6, diff_parts={K1 S2 I2 C1 F1}, diff_note={concepts: conservative field test | steps: curl condition; potential by integration; evaluate}]
Electric field intensity at a point $(x,y)$ is given by $\vec I=\left(12xy^3-4x\right)\hat i+18x^2y^2\hat j$ where all parameters are in S.I. Units. Then choose the correct statement(s) out of the following
\begin{qoptions}
\opt{A}{it represents a conservative electric field.}
\opt{B}{it represents a non-conservative electric field.}
\opt{C}{electric potential at point $(1\mathrm m,1\mathrm m)$ is $-4$V is electric potential at origin is zero.}
\opt{D}{electric potential at point $(1\mathrm m,1\mathrm m)$ cannot be defined if potential is taken zero at origin.}
\end{qoptions}
\end{bankq}
\begin{bankq}[id=TNAS03-P2-PHY-S1-Q03, source={TNAS - Test 03 (28 May 2023)}, paper=2, qno=3, type=multi, subtopic={Electric potential}, answer={A,C}, answer_src=key, difficulty=7, diff_parts={K1 S2 I2 C2 F1}, diff_note={concepts: charge crossing a potential step | steps: energy condition; angle of motion; fraction crossing}]
A positively charged particle is accelerated from the origin with the help of a voltage $4U$ (where $U>0$) to a certain velocity which lies on the $x-y$ plane under the angle of 30 degrees with respect to the $x$-axis (see figure). The electric potential $V(x,y)=V(x)$ only depends on the $x$-coordinate: if $3a<x<4a$ then $V\{x\}=U$ and if not then $V(x)=0$; besides the electrostatic force no other forces are applied to the particle and $a>0$.\\
For option C \& D consider the source of the charged particles is now a coaxial vacuum diode at the origin. Because of that the particles accelerated with the voltage $4U$ are moving isotropically (at equal quantities) to all directions of the $x-y$ plane; the $z$-directional component of the velocity for all particles is 0.
\bfig{92.5}{TNAS03_p2_phy_s1_q03_1.png}
\begin{qoptions}
\opt{A}{The trajectory of the charge on x-y plane will be \bfigi{31.8}{TNAS03_p2_phy_s1_q03_2.png}.}
\opt{B}{The trajectory of the charge on x-y plane will be \bfigi{36.2}{TNAS03_p2_phy_s1_q03_3.png}.}
\opt{C}{The part of all the emitted particles that reaches the region of space $x>4a$ is $\dfrac13$.}
\opt{D}{The part of all the emitted particles that reaches the region of space $x>4a$ is $\dfrac23$.}
\end{qoptions}
\end{bankq}
\begin{bankq}[id=TNAS03-P2-PHY-S3-Q01, source={TNAS - Test 03 (28 May 2023)}, paper=2, qno=1, type=numerical, subtopic={Conductors}, answer={60}, answer_src=key, difficulty=6, diff_parts={K1 S2 I2 C1 F1}, diff_note={concepts: tension between charged balls in a field | steps: induced charges; force on each; tension}]
Two metal balls of radius $R$ are connected to each other with a thin metal wire. The balls are in a homogeneous electric field of strength $E$, The length of the metal wire is $\ell$ ; furthermore $\ell>R$. The system is at equilibrium. Find the mechanical tension $T$ in the wire. If $T$ is given by $\alpha\times10^{-12}$ N fill $\alpha$. (Given $E=\sqrt{18}$ N/C, $\ell=3$ m, $R=2$ cm)
\end{bankq}
\begin{bankq}[id=TNAS04-P2-PHY-S1-Q01, source={TNAS - Test 04 (18 Jun 2023)}, paper=2, qno=1, type=multi, subtopic={Conductors}, answer={A,B,C}, answer_src=key, difficulty=6, diff_parts={K2 S2 I1 C1 F1}, diff_note={concepts: charged cavity in a conductor, induced charge and image force | steps: charge on the inner and outer surfaces; force on the cavity charge; force on the external charge; compare options}]
In the given figure `P' is a point charge of charge `q' placed inside on arbitrary shaped cavity of an uncharged conducting sphere of radius $R$ at a distance $R/2$ from centre `C', `S' is another point charge of charge ``2q'' placed at distance $2R$ from centre. PCS is a straight line. Choose this correct option(s) :-
\bfig{105.8}{TNAS04_p2_phy_s1_q01_1.png}
\begin{qoptions}
\opt{A}{P applies force on outer surface.}
\opt{B}{S applied force on inner surface.}
\opt{C}{Net electric force on P due to charge on outer surface is $\dfrac{2}{25}\dfrac{q^2}{\pi\epsilon_0R^2}$.}
\opt{D}{Net electric force on S due to ``sphere + P'' is rightward.}
\end{qoptions}
\end{bankq}
\begin{bankq}[id=TNAS05-P1-PHY-S1-Q03, source={TNAS - Test 05 (09 Jul 2023)}, paper=1, qno=3, type=single, subtopic={Electric potential}, answer={A}, answer_src=key, difficulty=7, diff_parts={K1 S3 I2 C2 F0}, diff_note={concepts: potential difference in a nonuniform charged shell | steps: field in the region by Gauss law; field outside; integrate field between R and 2R; work done; match}]
A spherical volume has charge density $\rho(r)=\dfrac{\rho r}{R}$ from $R$ to $2R$, where $r$ is distance from centre. There is no charge from $r=0$ to $r=R$ and $r=2R$ to infinity. Then work done to bring a charge $q$ slowly from $r=2R$ to $r=R$ is :-
\begin{qoptions}
\opt{A}{$\dfrac{11\rho R^2q}{24\varepsilon_0}$}
\opt{B}{$\dfrac{13\rho R^2q}{24\varepsilon_0}$}
\opt{C}{$\dfrac{13\rho R^2q}{12\varepsilon_0}$}
\opt{D}{$\dfrac{11\rho R^2q}{12\varepsilon_0}$}
\end{qoptions}
\end{bankq}
\begin{bankq}[id=TNAS05-P1-PHY-S1-Q05, source={TNAS - Test 05 (09 Jul 2023)}, paper=1, qno=5, type=single, subtopic={Gauss's law}, answer={D}, answer_src=key, difficulty=7, diff_parts={K1 S3 I2 C2 F0}, diff_note={concepts: flux through a triangular plate from a line charge | steps: field of the wire at a distance; geometry of the plate; flux element; integrate; evaluate}]
An infinite long wire having linear charge density $\lambda$ is placed parallel to $x$-axis as shown. Find flux passing through the triangular plate placed in $x$-$z$ plane?
\bfig{145.2}{TNAS05_p1_phy_s1_q05_1.png}
\begin{qoptions}
\opt{A}{$\dfrac{\lambda a}{8\varepsilon_0}$}
\opt{B}{$\dfrac{\lambda a}{12\varepsilon_0}$}
\opt{C}{$\dfrac{\lambda a}{16\varepsilon_0}$}
\opt{D}{$\dfrac{\lambda a}{4\varepsilon_0}$}
\end{qoptions}
\end{bankq}
\end{document}
